% concatenated from main.tex

%--------------- Personalize your document here ---------------


\documentclass[a4paper,12pt]{article}

\title{Orthospectrum and simple orthospectrum rigidity: finiteness and genericity.} %Enter the title 
\date{2024} % insert a specific date	

\author{Nolwenn Le Quellec} % Enter your name

\newcommand{\address}{Univ Gustave Eiffel, Univ Paris Est Creteil, CNRS, LAMA UMR8050 F-77447 Marne-la-Vallée, France}
\newcommand{\email}{nolwenn.le-quellec@univ-eiffel.fr}



\usepackage[left=25mm,top=30mm,right=25mm,bottom=30
mm]{geometry}
\usepackage{booktabs}
\usepackage[usestackEOL]{stackengine}
\usepackage[T1]{fontenc}
\usepackage[sfdefault]{atkinson}
\usepackage[utf8]{inputenc}
\usepackage{bm}
\usepackage{graphicx}
\usepackage{subcaption}
\usepackage[all]{xy}
\usepackage{amsthm}
\usepackage{amsfonts}
\usepackage{mathtools}
\usepackage{animate}
\usepackage{faktor}
\usepackage[square,numbers,sort]{natbib}
\usepackage[french, english]{babel}
\usepackage{xcolor}
\usepackage{float}
\usepackage{enumitem,kantlipsum}
\usepackage{amssymb}
\usepackage{listings}
\usepackage{stmaryrd}

\definecolor{PrettyBlue}{RGB}{11, 83, 148}
\usepackage[breaklinks,colorlinks=true,allcolors=PrettyBlue]{hyperref}


\linespread{1}




\newtheorem{theorem}{Theorem}[section]
\newtheorem{proposition}[theorem]{Proposition}
\newtheorem{lemme}[theorem]{Lemma}
\newtheorem{corollaire}[theorem]{Corollary}

\newtheorem*{mainTheoremAa}{\theoremnameaa}
\newcommand{\theoremnameaa}{\hypertarget{thm:Main1aIntro}{Theorem~\ref{thm:OrthoFiniteb}}}
\newenvironment{Aa}[1]{\renewcommand{\theoremnameaa}{#1}\begin{mainTheoremAa}}{\end{mainTheoremAa}}

\newtheorem*{mainTheoremAb}{\theoremnameab}
\newcommand{\theoremnameab}{\hypertarget{thm:Main1bIntro}{Theorem~\ref{thm:OrthoFiniteSimple}}}
\newenvironment{Ab}[1]{\renewcommand{\theoremnameab}{#1}\begin{mainTheoremAb}}{\end{mainTheoremAb}}

\newtheorem*{mainTheoremB}{\theoremnameb}
\newcommand{\theoremnameb}{\hypertarget{thm:Main2Intro}{Theorem~\ref{thm:Main2}}}
\newenvironment{B}[1]{\renewcommand{\theoremnameb}{#1}\begin{mainTheoremB}}{\end{mainTheoremB}}

\newtheorem*{mainTheoremC}{\theoremnamec}
\newcommand{\theoremnamec}{\hypertarget{thm:Main3Intro}{Theorem~\ref{thm:Main3}}}
\newenvironment{C}[1]{\renewcommand{\theoremnamec}{#1}\begin{mainTheoremC}}{\end{mainTheoremC}}

\theoremstyle{definition}
\newtheorem{definition}[theorem]{Definition}
\newtheorem{ex}[theorem]{Example}
\newtheorem*{Ack}{Acknowledgments}
\newtheorem*{Abs}{Abstract}
\newtheorem*{stm}{AI Statement}
\theoremstyle{remark}
\newtheorem{rem}[theorem]{Remark}



\begin{document}

\maketitle

\pagenumbering{gobble}% Remove page numbers (and reset to 1)



\begin{Abs}
    We study the orthospectrum and the simple orthospectrum of compact hyperbolic surfaces with geodesic boundary. We show that there are finitely many hyperbolic surfaces sharing the same simple orthospectrum and finitely many hyperbolic surfaces sharing the same orthospectrum. Then, we show that generic surfaces are determined by their orthospectrum and by their simple orthospectrum. We conclude with the example of the one-holed torus which is determined by its simple orthospectrum.
\end{Abs}

\begin{stm}
    This work was completed without use of artificial intelligence. The author does not consent for all or part of this work to be used for training of artificial intelligence models.
\end{stm}

%\tableofcontents


\pagenumbering{arabic}% Arabic page numbers (and reset to 1)

% This is how you can organize your document
\section{Introduction}

In a by now famous paper \cite{kac}, Kac asked "Can one hear the shape of a drum?". The question can be formalized mathematically using the fact that the sound made by a drum is linked to the frequencies at which a drumhead vibrates, that is, the spectrum of the Laplacian. Kac's question then becomes: "Does the spectrum of the Laplacian of a planar domain determine the domain itself?". This question was then also asked for general Riemannian manifolds. In the case of closed hyperbolic surfaces, Huber and Selberg (see \cite[Chapter~7.1]{buser}) showed that the Laplacian spectrum determines and is determined by the length spectrum, which is the set of lengths of closed geodesics counted with multiplicities. This shifted the question to "Does the length spectrum of a hyperbolic surface determine the metric up to isometry?". 
In 1978, Vignéras provided the first example of isospectral non-isometric closed hyperbolic surfaces \cite{vigneras} showing that the answer to the question is negative, and in 1985, Sunada gave a general criterion to construct isospectral non-isometric manifolds \cite{sunada}. On the other hand, McKean proved that there is a finite number of isospectral non-isometric hyperbolic surfaces \cite{McKean}. In 1979, Wolpert showed that a generic hyperbolic surface is determined by its length spectrum \cite{Wolpert}. In 1985, Haas proved that the one-holed hyperbolic torus with a fixed boundary length is determined by the length spectrum \cite{Haas}, and later Buser and Semmler removed the condition on the boundary length \cite{BuSem}.
In the case of the simple spectrum (where only the simple closed geodesics are considered) the question is still open: there are no known examples of non-isometric hyperbolic surfaces with the same simple length spectrum. There is, however, a result by Baik, Choi and Kim showing that generic hyperbolic surfaces are also determined by their simple length spectrum \cite{BaikChoiKim}. 

In 1993, Basmajian introduced the orthospectrum of hyperbolic surfaces with boundary, defined as the set of lengths of geodesic arcs orthogonal to the boundary (called orthogeodesics) counted with multiplicities \cite{Basmajian}. This object, analogous to the length spectrum, aims at taking the boundary of a surface more into account. In particular, Basmajian gave a formula to compute the boundary length of a hyperbolic surface from its orthospectrum \cite{Basmajian}. Motivated by Kac's question, Masai and McShane considered the analogous problem for the orthospectrum: "Does the orthospectrum of a hyperbolic surface determine the metric up to isometry?". They showed that there is a finite number of non-isometric hyperbolic surfaces with one boundary component sharing the same orthospectrum \cite{orthoSys}. In the same paper, they showed that the one-holed torus is determined by its orthospectrum. They also gave a general construction for isospectral non-isometric hyperbolic surfaces, thus answering the question in the negative.\\

In this paper, we investigate further Masai and McShane's question, both for the orthospectrum and the simple orthospectrum. With the simple orthospectrum define as the multiset of lengths of simple orthogeodesics, counted with multiplicities. Note that even if the orthospectrum and simple orthospectrum seem related, there is no known way to deduce one from the other. Denoting by $\mathcal{O}(X)$ and $\mathcal{O}_S(X)$ the orthospectrum and the simple orthospectrum of a hyperbolic surface $X$, our main results are the following.


\begin{mainTheoremAa}\label{thm:Main1aIntro}
    Let $S_g^b$ be a compact, genus~$g$ surface of negative Euler characteristic with~$b$ boundary components. Let $X$ be a hyperbolic structure on~$S_g^b$ with geodesic boundary. Then, up to isometry, there is a finite number of hyperbolic structures $Y$ on~$S_g^b$ such that
    \[
        \mathcal{O}(Y)=\mathcal{O}(X).
    \]
\end{mainTheoremAa}

This theorem extends Masai and McShane's result from surfaces with a single boundary component to surfaces with several boundary component. The proof is based on Masai and McShane's proof. The following theorem states a variation of this result for the simple orthospectrum.

\begin{mainTheoremAb}\label{thm:Main1bIntro}
    Let $S_g^b$ be a compact, genus~$g$ surface of negative Euler characteristic with~$b$ boundary components. Let $X$ be a hyperbolic structure on~$S_g^b$ with geodesic boundary. Then, up to isometry, there is a finite number of hyperbolic structures $Y$ on~$S_g^b$ such that
    \[
        \mathcal{O}_S(Y)=\mathcal{O}_S(X).
    \]
\end{mainTheoremAb}

The proof of this theorem is inspired by the compactness argument in Masai and Mcshane's proof, an argument that dates back to Wolpert, but the way we show that the set of hyperbolic surfaces with the same simple orthospectrum is compact is completely different. Instead of relying on orthogeodesics with one self-intersection as Masai and McShane do, our proof uses arc and curve counting results by Bell and Mirzahani. \\

We also establish an analogue of Wolpert's result both for the orthospectrum and the simple orthospectrum.



%%%%%%%%%%%%

\begin{mainTheoremB}\label{thm:Main2Intro}
    Generic surfaces in $\mathrm{Teich}(S_g^b)$ are determined, up to isometry, by their orthospectrum. Similarly, generic surfaces in $\mathrm{Teich}(S_g^b)$ are also determined, up to isometry, by their simple orthospectrum. 
\end{mainTheoremB}

We were inspired by Wolpert's proof of the generic determination by the length spectrum. Main differences between the proofs lie in the preparation and the use of a different type of coordinates for the Teichmüller space in the main part of the proof. Finally, we show:

\begin{mainTheoremC}\label{thm:Main3Intro}
    Let $T$ and $T'$ be two hyperbolic structures with geodesic boundary on the one-holed torus.
    Then $T$ and $T'$ are isometric if and only if $\mathcal{O}_S(T)=\mathcal{O}_S(T')$.
\end{mainTheoremC}

This paper is organized as follows. We recall in Section~\ref{sec:1} properties of hyperbolic surfaces and geodesics that will be needed in the rest of the paper. In Section \ref{sec:2}, we prove Theorem~\ref{thm:OrthoFiniteb} and Theorem~\ref{thm:OrthoFiniteSimple}. We prove Theorem~\ref{thm:Main2} in Section~\ref{sec:3}. Finally, in Section~\ref{sec:4} we show Theorem~\ref{thm:Main3}.

%%%%%%%%%%%%%%%%%%%%%%%%%%%%%%%%

\begin{Ack}
    We appreciate the support, guidance and help of Federica Fanoni and Stéphane Sabourau, both PhD advisors of the author, in particular for their help in the writing. \\
    
    The author would also like to thank Nhat Minh Doan for finding a mistake in the first version of the paper and Marie Trin for her help to understand curve counting results.
\end{Ack}










\section{Preliminaries}\label{sec:1}


\subsection{Curves and arcs on surfaces.}
Throughout this paper, $S_g^b$ will denote an orientable, compact, genus~$g$ surface of negative Euler characteristic with $b>0$ boundary components. Moreover, arcs are always defined with endpoints on the boundary and closed curves are always considered primitive.


\begin{definition}
    A closed curve on $S_g^b$ is \emph{essential} if it is not homotopic to a boundary component or to a point. It is \emph{simple} if it has no self-intersection.
\end{definition}


\begin{definition}
    A \emph{multi-curve} is a finite formal sum $\gamma = \sum_i\alpha_i\gamma_i$ of pairwise disjoint, non-homotopic, essential, primitive, simple closed curves $\gamma_i$ with positive coefficients $\alpha_i$ known as weights. It is \emph{integral} if the weights are integers and \emph{rational} if the weights are rational numbers.\\
    The set of integral multi-curves on $S_g^{b}$ is denoted
    \[
        \mathcal{ML}_{g,b}(\mathbb{Z}).
    \]
\end{definition}
Note that a simple closed curve can be seen as an integral multi-curve with only one weight that is equal to $1$.
This notation of the set of integral multi-curves comes from the notation of the set of laminations that contains it. See~\cite{thurstonWork} and~\cite{VivekaJuan} for more details on laminations. 



\begin{definition}
    A \emph{pair of pants} is a surface homeomorphic to $S_0^3$.
    A \emph{pants decomposition} of $S_g^b$ is a maximal collection of pairwise non-homotopic, disjoint, essential simple closed curves on $S_g^b$.
\end{definition}

\begin{rem}
    The cardinality of a pants decomposition is $3g+b-3$ and a pants decomposition cuts $S_g^b$ into $2g+b-2$ pairs of pants \cite{TeichBord}. 
\end{rem}


%%%%%%%%

When we study surfaces with boundary, it can be useful to change the perspective from the closed curve/closed geodesic point of view to the arc/orthogeodesic point of view. 


\begin{definition}
    An arc on $S_g^b$ is \emph{essential} if it is not homotopic relatively to the boundary, into a boundary component. It is \emph{simple} if it has no self-intersection.
    
    An \emph{orthogeodesic} of a compact hyperbolic surface is the shortest geodesic representative of the homotopy class relative to the boundary of an essential arc.
\end{definition}

%%%%%%%%%
\begin{rem}
    Endpoints of an orthogeodesic are orthogonals to the boundary \cite{Basmajian}.
\end{rem}
%%%%%%%%%%%

\begin{definition}
    A \emph{hexagon decomposition} of a surface $S_g^b$ is a maximal collection of pairwise non-homotopic and disjoint simple essential arcs on $S_g^b$.
\end{definition}

\begin{rem}
    There always exists a hexagon decomposition on a negative Euler characteristic surface.
    The cardinality of a hexagon decomposition is $6g+3b-6$ and the hexagon decomposition cuts $S_g^b$ into $4g+2b-4$ hexagons. See \cite{Akira}.
\end{rem}


When we endow $S_g^b$ with a hyperbolic metric, and cut the surface along the orthogeodesic representatives of a hexagon decomposition, the surface decomposes into right-angled hexagons.
\begin{figure}[H]
    \centering
    \includegraphics[height=10cm]{figures/HexaToreTroue.png}
    \caption{Example of a hexagon decomposition of the one-holed torus.}
    \label{fig:ExHexa}
\end{figure}


There is a link between pants and hexagon decomposition through the concept of double of surfaces.

\begin{definition}
    Let $S$ be a surface of genus $g$ with $b>0$ boundary components. The \emph{double of $S$} is the surface obtained by gluing two copies of $S$ along their corresponding boundary components. 
    When $S$ is endowed with a hyperbolic metric with geodesic boundary components, the double of $S$ is endowed with the hyperbolic metric which coincides with the hyperbolic metric of $S$ on each copy of $S$ and such that, the reflection~$R$ between the two copies of $S$ along their boundary is an isometry. \\
    Let $a$ be an arc of $S$. We call the \emph{double of $a$} the union of the two copies of $a$ on the double of $S$.
\end{definition}


\begin{rem}
    The double of $S_g^b$ is the surface $S_{2g+b-1}$ of genus $2g+b-1$ with no boundary.
\end{rem}

\begin{lemme}\label{lem:HexaToPant}
    Let $\mathcal{H}$ be a hexagon decomposition of $S_g^b$. Denote by $S_{2g+b-1}$ the double of $S_g^b$ and by $R$ the reflection associated to it.
    Then, $\mathcal{H} \cup R(\mathcal{H})$ is a pants decomposition on $S_{2g+b-1}$.
\end{lemme}

\begin{proof}
    The double of every arc of a hexagon decomposition forms a simple closed curve of $S_{2g+b-1}$. By construction, $\mathcal{H} \cup R(\mathcal{H})$ is a set of $6g+3b-6$ pairwise non-homotopic, disjoint, essential, simple, closed curves on $S_{2g+b-1}$. Therefore, it is a maximal collection of such curves, and hence a pants decomposition.
\end{proof}

Now, let us define the geometric intersection number between closed curves and arcs.

\begin{definition}
    Let $\alpha$ and $\beta$ be closed curves or arcs on a surface $S$. The \emph{geometric intersection number} $i(\alpha,\beta)$ between $\alpha$ and $\beta$ is the minimum number of intersection counted with multiplicities between curves/arcs in their homotopy class (relative to the boundary).
\end{definition}

Note that, if $\alpha$ and $\beta$ are distinct closed geodesics or orthogeodesics \mbox{$i(\alpha, \beta)=|\alpha \cap \beta |$}.

We can define coordinates in the space of multi-curve that topologically identified each integral multi-curve. 

\begin{definition}
    Let $\mathcal{P}=\{\alpha_1,...,\alpha_{3g-3+b}\}$ be a pants decomposition on $S_g^{b}$. For all multi-curve $\gamma$ we define its Dehn-Thurston coordinates as $(m_i,t_i)_{i=1}^{3g+b-3}$ with $m_i$ the intersection number between $\gamma$ and $\alpha_i$, and $t_i$ the twisting number of $\gamma$ around $\alpha_i$ as defined in~\cite{DehnThurstonCoord}.
\end{definition}


Let $Z(\mathcal{P})$ be the set of $(m_i,t_i)^{3g-3+b}_{i=1} \in (\mathbb{Z}_+ \times \mathbb{Z})^{3g-3+b}$ satisfying the following conditions:
\begin{itemize}
    \item If $m_i =0$, then $t_i\geqslant0$,
    \item If $\alpha_{i_1}$, $\alpha_{i_2}$ and $\alpha_{i_3}$ are the boundary components of a pair of pants in $S_g^{b}$, then $m_{i_1}+m_{i_2}+m_{i_3}$ is even.
\end{itemize}
Then we have the Dehn theorem that guarantees that Dehn-Thurston coordinates uniquely determine integral multi-curves~\cite{PennerHarer}.

\begin{theorem}\label{thm:DehnThurston}
    For any pair of pants decomposition $\mathcal{P}$ on $S_g^{b}$, the map
    \[
        DT : \mathcal{ML}_{g,b}(\mathbb{Z}) \to Z(\mathcal{P})
    \]
    is a bijection.
\end{theorem}




\subsection{Teichmüller space and moduli space.}


The hyperbolic surfaces studied in this paper either live in the Teichmüller space or the moduli space of $S_g^b$. Moreover, we assume they all have a geodesic boundary. 

\begin{definition}
    The \emph{Teichmüller space} $\mathrm{Teich}(S_g^b)$ of $S_g^b$ is the set of homotopy classes of hyperbolic structures on $S_g^b$.
\end{definition}

\begin{definition}
    The \emph{mapping class group} $\mathcal{MCG}(S_g^b)$ of $S_g^b$ is the quotient of the group of orientation-preserving homeomorphisms of $S_g^b$ by the subgroup of homeomorphisms isotopic to the identity.
    \[
        \mathcal{MCG}(S_g^b) = \mathrm{Homeo}^+(S_g^b)/\mathrm{Homeo}_0(S_g^b).
    \]
\end{definition}

\begin{definition}
    The \emph{moduli space} $ \mathcal{M}(S_g^b)$ is the space of isometry classes of hyperbolic surfaces homeomorphic to $S_g^b$.
\end{definition}

The moduli space is the quotient space
\[
    \mathcal{M}(S_g^b) = \mathrm{Teich}(S_g^b)/\mathcal{MCG}(S_g^b)
\]
and Teichmüller space is the universal cover of $\mathcal{M}(S_g^b)$.\\


The following function on $\mathrm{Teich}(S_g^b)$ will be used throughout this paper.
\begin{definition}
    For any essential closed curve/arc $\alpha$ on $S_g^b$ and any hyperbolic surface $\chi \in \mathrm{Teich}(S_g^b)$, the \emph{length function} $\ell_{\chi}(\alpha)$ is the length of the shortest curve/arc in the homotopy class of $\alpha$ (relative to the boundary if $\alpha$ is an arc) on~$\chi$. 
    If $\mathcal{C}$ is a collection of curves and arcs, we define $\ell_{\chi}(\mathcal{C}) = \sum_{\alpha \in \mathcal{C}} \ell_{\chi}(\alpha) $ and if $\gamma=\sum_i^k a_i\gamma_i$ is a multi-curve, we define $\ell_\chi(\gamma)=\sum_i^k a_i\ell_\chi(\gamma_i)$.
\end{definition}

We may write $\ell_X(\alpha) $ with $X \in \mathcal{M}(S_g^b)$ instead of its lift $\chi \in \mathrm{Teich}(S_g^b)$ when there is no ambiguity on the marking of $\chi$. When there is no ambiguity about the surface, we may just write $\ell(\alpha)$. \\

With the Dehn-Thurston coordinates we define another kind of length with respect to a pants decomposition: the combinatorial length on integral multi-curves.
\begin{definition}
    Let $\mathcal{P}=\{\alpha_1,...,\alpha_{3g-3+b}\}$ be a pants decomposition on $S_g^{b}$ and $\gamma$ an integral multi-curve on $S_g^{b}$ of Dehn-Thurston coordinates $(m_i,t_i)^{3g-3+b}_{i=1}$. Its combinatorial length on $\chi$ is
    \[
        \mathcal{L}_\mathcal{P}(\chi,\gamma) = \sum_{i=1}^{3g-3+b} \left( m_i S(\ell_\chi(\alpha_i)) + |t_i|\ell_\chi(\alpha_i)  \right)
    \]
    with $S(x)=\sinh^{-1}\left( \frac{1}{\sinh(x/2)} \right)$.
\end{definition}

Let us define the twist function that depends on a pants decomposition fixed on $S$ and a system of curves defined relatively to the pants decomposition as follow. 
We fix a pants decomposition $\mathcal{P}$ on $S$. Let $\mathcal{P}'$ be a collection of disjoint simple closed curves on $S$ such that for each pair of pants $P$ in $S \setminus \mathcal{P}$, $\mathcal{P}' \cap P$ consist of $3$ arcs each connecting a different pair of boundary components of $P$. Such a system of curves always exists regardless of the choice of $\mathcal{P}$ and we always fix one along fixing a pants decomposition. 


\begin{definition}
    For any fixed pants decomposition $\mathcal{P}$, any curve $\alpha \in \mathcal{P}$ and any hyperbolic structure $\chi = (X,\varphi) \in \mathrm{Teich}(S)$, the \emph{twist function} $t_{\chi}(\alpha)$ is defined as follows. \\
    
    Let $\alpha \in \mathcal{P}$, we give an orientation to $\alpha$ and denote $P_l$ and $P_r$ the pair of pants in $S \setminus \mathcal{P}$ respectively on the left and on the right of $\alpha$ (they may be the same pair of pants). Let $\eta$ be one of the two components of $(P_l \cup P_r \cup \alpha) \cap \mathcal{P}'$ that intersect $\alpha$.\\
    
    Let $\delta$ be the geodesic representative of $\varphi(\alpha)$ on $X$. Let $\tau_r$ and $\tau_l$ be the orthogeodesics in $\varphi(P_r)$ and $\varphi(P_l)$ between $\delta$ and the boundary component hit by $\varphi(\eta)$. Denote by $p_r$ and $p_l$ the points on $\delta$ hit by $\tau_r$ and $\tau_l$. Relatively to the boundary of $\varphi ( (P_r \cup P_l \cup \alpha) )$, the arc $\varphi(\eta)$ is homotopic to $\tau_r . \delta^k . \tau_l$ for some $k \in \mathbb{Z}$. We define the \emph{twist function} as 
    \[
        t_{\chi}(\alpha)=k + \frac{d(p_r,p_l)}{\ell_\chi(\alpha)}
    \]
    with $d(p_r,p_l)$ the signed distance between $p_r$ and $p_l$ along the orientation of $\varphi(\alpha)$.    
\end{definition}

A well-known system of coordinates on the Teichmüller space are the Fenchel-Nielsen coordinates defined with the length and twist function relatively to a pants decomposition. For a fixed pants decomposition we can then relate the combinatorial length to the length function through the following theorem by Mirzakhani~\cite[Proposition~3.5]{mirzakhani}. 
\begin{theorem}
    Let $\mathcal{P}$ be a pants decomposition on $\chi \in \mathrm{Teich}(S_g^b)$ such that for all $\alpha \in \mathcal{P}$, we have $\ell_\chi(\alpha) \leqslant L$ and $t_\chi(\alpha) \in [0,1]$. Then for any integral multi-curve on $\chi$, we have
    \begin{align}
        \frac{1}{c(L)}\mathcal{L}_\mathcal{P}(\chi,\gamma) \leqslant \ell_\chi(\gamma) \leqslant c(L) \mathcal{L}_\mathcal{P}(\chi,\gamma),\label{in:combiLength}
    \end{align}
    with $c$ that only depends on $L, g$ and $b$.
\end{theorem}

Let us define another system of coordinates with only the length function and a hexagon decomposition. A right-angled hexagon is determined up to isometry by the lengths of three pairwise non-consecutive sides \cite[Theorem 3.5.14]{Rat}. Thus, if we fix a hexagon decomposition on $S_g^b$, the lengths of the orthogeodesic representatives of the hexagon decomposition on $S_g^b$ determine a hyperbolic surface in the Teichmüller space. 
Furthermore, Ushijima showed in \cite[Theorem 4.1]{Akira} the following theorem:

\begin{theorem}\label{HexaCoord}
    Given a hexagon decomposition $\mathcal{H}=\{\alpha_1,...,\alpha_{6g+3b-6}\}$ on $S_g^b$, the map
    \[
        \varphi_{\mathcal{H}} : \mathrm{Teich}(S_g^b) \to  \mathbb{R}^{6g+3b-6}_+
    \]
    defined by
    \[
        \varphi_{\mathcal{H}}(\chi)=(\ell_{\chi}(\alpha_1),...,\ell_{\chi}(\alpha_{6g+3b-6})) 
    \]
    is a homeomorphism.
\end{theorem}



\begin{definition}\label{def:omegaSurface}
    For a fixed hexagon decomposition $ \mathcal{H} = \{\alpha_1,...,\alpha_{6g+3b-6}\}$, we call \\
    $\varphi_{\mathcal{H}}(\centerdot)=(\ell_{\centerdot}(\alpha_1),...,\ell_{\centerdot}(\alpha_{6g+3b-6}))$ \textit{Ushijima coordinates} function and denote by 
    \[
      \chi^{\omega}=\varphi_{\mathcal{H}}^{-1}(\omega)   
    \] 
    the surface in $\mathrm{Teich}(S_g^b)$ associated to $\omega \in \mathbb{R}_+^{6g+3b-6}$.
\end{definition}

%%%%%%%%%%%%
The Teichmüller space of $S_g^b$ admits a real analytic structure and the Fenchel-Nielsen coordinates are real analytic (independently from the pants decomposition)~\cite{TeichBord}. By identifying $\mathrm{Teich}(S_g^b)$ with a real analytic subvariety of $\mathrm{Teich}(S_{2g+b-1})$, we obtain that this structure is compatible with the one we get via Ushijima coordinates: fix a hexagon decomposition $\mathcal{H}=\{ \alpha_1,...,\alpha_{6g+3b-6} \}$ on $S_g^b$, by Lemma~\ref{lem:HexaToPant} this hexagon decomposition induces a pants decomposition on $S_{2g+b-1}$, the double of $S_g^b$, we denote it $\mathcal{P}=\{ \alpha_1',...,\alpha_{6g+3b-6}' \}$ where $\alpha_i'$ is the double of $\alpha_i$. We denote by $T(S_g^b)$ the subset of surfaces in $\mathrm{Teich}(S_{2g+b-1})$ that are double of surfaces in $\mathrm{Teich}(S_g^b)$, which is identified with $\mathrm{Teich}(S_g^b)$. We have 
    \[
        T(S_g^b)=\{ \chi \in \mathrm{Teich}(S_{2g+b-1}) \mid t_{\alpha'_i}(\chi)=0, 1 \leqslant i \leqslant 6g+3b-6 \}
    \]
with $t_{\alpha'_i}(.)$ the twists parameters in the Fenchel-Nielsen coordinates. The length parameters are real analytic on $T(S_g^b)$ and by identifying $\mathrm{Teich}(S_g^b)$ with $T(S_g^b)$ Ushijima coordinates $(\ell(\alpha_1),...,\ell(\alpha_{6g+3b-6}))=(\frac{\ell(\alpha_1')}{2},...,\frac{\ell(\alpha_{6g+3b-6}')}{2})$ are real analytic. This does not depend on the hexagon decomposition. Through the same process with a pants decomposition on $S_{2g+b-1}$ induced by a pants decomposition on $S_g^b$, we show that Fenchel-Nielsen coordinates are real analytic for the same analytic structure as Ushijima coordinates.

\begin{lemme}\label{lem:RealAnalLength}
    For any arc $\beta$, the function 
    \begin{align*}
        \ell_{\centerdot} (\beta) : \mathrm{Teich}(S_g^b) & \to \mathbb{R}^+ \\
        \chi & \mapsto \ell_{\chi}(\beta)
    \end{align*}
     is real analytic on $\mathrm{Teich}(S_g^b)$.
\end{lemme}

\begin{proof}
    Let $\beta$ be an arc on $S_g^b$ and $\beta'$ be the double of $\beta$ on $S_{2g+b-1}$. By \cite[Lemma~10.2.3]{buser} (see also \cite[Proposition~2.4]{McShaneParlier}), we know that for any closed curve $\beta'$ the function $\ell_{\centerdot}(\beta')$ is real analytic on $\mathrm{Teich}(S_{2g+b-1})$ and thus on $T(S_g^b)$. By identifying $T(S_g^b)$ and $\mathrm{Teich}(S_g^b)$, we get that $\ell_{\centerdot} (\beta)=\frac{\ell_{\centerdot} (\beta')}{2} $ is real analytic on $\mathrm{Teich}(S_g^b)$.
\end{proof}

%%%%%%%%%%%%%%


%%%%%%

\begin{definition}
    Given two surfaces $\chi$, $\Upsilon \in \mathrm{Teich}(S_g^b)$ and $K \geqslant 1$, a homeomorphism $\phi : \chi \to \Upsilon$ is said to be \emph{$K$-quasiconformal} if its distributional derivatives are locally in~$L_2$ and
    \[
        |\phi_{\overline{z}}|\leqslant \frac{K-1}{K+1} |\phi_{z}| \text{ a.e.}
    \]
\end{definition}

For any $\chi, \Upsilon \in \mathrm{Teich}(S_g^b)$, there exists a unique $K$-quasiconformal map $\phi_{\chi, \Upsilon}$ between~$\chi$ and $\Upsilon$ with $K$ minimal (see \cite{TeichBord} for more details). The Teichmüller metric~$d_{\mathrm{Teich}}$ is defined by $d_{\mathrm{Teich}}(\chi,\Upsilon)=\frac{1}{2}\log(K)$; see \cite{teich}. \\

%%%%%%%

Now, let us state an extension of Wolpert's lemma~\cite[12.3.2]{teich} to any closed curve, simple or not, due to Buser~\cite[Theorem~6.4.3]{buser}.

\begin{lemme}[Wolpert's lemma]\label{lemme:Wolpert}
    Let $\phi : \chi_1 \to \chi_2$ be a $K$-quasiconformal homeomorphism between two hyperbolic surfaces $\chi_1$ and $\chi_2$. For any isotopy class $c$ of a closed curve in $\chi_1$, the following inequalities hold:
    \[
        \frac{\ell_{\chi_1}(c)}{K} \leqslant \ell_{\chi_2}(\phi(c)) \leqslant K \ell_{\chi_1}(c).
    \]
\end{lemme}

The previous lemma applied to the double of the surface implies the following.

\begin{lemme}[Orthogeodesic Wolpert's lemma]\label{lemme:OrthoWolpert}
    Let $\phi : \chi_1 \to \chi_2$ be a $K$-quasiconformal homeomorphism between two hyperbolic surfaces $\chi_1$ and $\chi_2$. For any isotopy class $c$ of arc in $\chi_1$, the following inequalities hold:
    \[
        \frac{\ell_{\chi_1}(c)}{K} \leqslant \ell_{\chi_2}(\phi(c)) \leqslant K \ell_{\chi_1}(c).
    \]
\end{lemme}


By definition of the Teichmüller metric, for any $\chi_1, \chi_2 \in \mathrm{Teich}(S_g^b)$ with \\$d_{\mathrm{Teich}}(\chi_1,\chi_2) = \log(K)/2$, there is a $K-$quasiconformal homeomorphism between $\chi_1$ and $\chi_2$. By Lemma~\ref{lemme:OrthoWolpert}, for any $\chi_1, \chi_2 \in \mathrm{Teich}(S_g^b)$ with $d_{\mathrm{Teich}}(\chi_1,\chi_2) \leqslant \log(K)/2$ and any isotopy class $c$ of simple arcs on $S_g^b$, we have $ \frac{\ell_{\chi_1}(c)}{K} \leqslant \ell_{\chi_2}(c) \leqslant  K\ell_{\chi_1}(c)$.\\

%%%%%%%%%%

Then, we state a corollary of Lemma~\ref{lemme:OrthoWolpert}.

\begin{corollaire}\label{cor:Wolp}
    Let $\mathcal{Q} \subset \mathrm{Teich}(S_g^b)$ be a compact subset. There exists a constant $q \geqslant 1$ which depends only on $\mathcal{Q}$ such that
    \[
        \frac{\ell_{\chi} (\beta)}{q} \leqslant \ell_{\chi'}( \beta) \leqslant q \ell_{\chi} (\beta)
    \]
    for any $\chi, \chi' \in \mathcal{Q}$ and any orthogeodesic $\beta$.
\end{corollaire}

\begin{proof}
    Let $q>0$ such that $\text{diam } \mathcal{Q}= \frac{1}{2} \log(q)$. Then for any $\chi, \chi' \in \mathcal{Q}$, \\
    we have $d_{\mathrm{Teich}}(\chi,\chi') \leqslant \log(q)/2$. By Lemma~\ref{lemme:OrthoWolpert}, we have
    \[
        \frac{\ell_{\chi} (\beta)}{q}  \leqslant \ell_{\chi'} (\beta) \leqslant q \ell_{\chi} (\beta).
    \]
    
\end{proof}

%%%%%%%%%%%%%%%%%

Thanks to the length function we can define the two following quantities that will regularly appear throughout this paper.

\begin{definition}
    The \emph{systole} $\mathrm{sys}(X)$ of a hyperbolic surface $X$ is the shortest length of an essential closed curve on $X$. Note that the systole is realized by the length of a simple closed geodesic, unless $X$ is a pair of pants.
\end{definition}

Similarly to the systole, we also define the orthosystole.

\begin{definition}
    The \emph{orthosystole} $\mathrm{osys}(X)$ of a hyperbolic surface $X$ is the shortest length of an orthogeodesic on $X$.
\end{definition}

Moreover, let us define the injectivity radius:

\begin{definition}
    The \emph{injectivity radius of a hyperbolic surface $X$ at $p \in X$}, denoted $r_p(X)$, is defined as the supremum of all $r$ for which $U_p^r=\{ q \in X \mid \mbox{dist}(p,q) < r \}$ is diffeomorphic to a disk or half a disk. The injectivity radius of $X$ is defined as\\ $r_{\mbox{inj}}(X)=\inf\{ r_p(X) \mid p \in X \}$.
\end{definition}

This quantity is related to the systole as by Lemma~\cite[4.8 chap.I]{Hummel} we have 
\[
    r_{\mbox{inj}}(X)= \frac{\min\{ \mathrm{sys}(X) , \min_{0 \leqslant j \leqslant b-1}\{ \ell_X(\gamma_j) \} \}}{2}, 
\]
where the $\gamma_j$ are the boundary components of $X$.\\

The injectivity radius allows us to define the thick-part of a hyperbolic surface.
\begin{definition}
    The \emph{$\varepsilon$-thick part} of a hyperbolic surface $X$ is defined as
    \[
        thick_\varepsilon(X) = \{ p \in X \mid r_p(X) \geqslant \varepsilon  \}.
    \]
\end{definition}


%%%%%%%%%%%%%%%%%%%%%%%%%%%%%%%%%%%%%%%

Finally, let us define an interesting subset of the moduli space.

\begin{definition}
    Let $\gamma_0,...,\gamma_{b-1}$, be the boundary components of $S_g^b$. We define the set
    \[
        \mathcal{M}_{A, \varepsilon}(S_g^b) = \{X \in \mathcal{M}(S_g^b) \mid \mathrm{sys}(X) \geqslant \varepsilon \text{ and } A \geqslant \ell_X(\gamma_0),...,\ell_X(\gamma_{b-1}) \geqslant \varepsilon \}
    \]
\end{definition}

The following result by Parlier \cite{bers} will help us define a property of $\mathcal{M}_{A, \varepsilon}(S_g^b)$.

\begin{theorem}\label{theorem:Bers}
    Let $X$ be a finite area hyperbolic surface, possibly with geodesic boundary~$\partial X$. Then $X$ admits a pants decomposition where each curve is of length at most
    \[
        B = \max\{ \ell(\partial X), \mathrm{area}(X) \}.    
    \]
     
\end{theorem}

\begin{theorem}\label{theorem:mumford}
    For all $ A \geqslant \varepsilon > 0$, $\mathcal{M}_{A, \varepsilon}(S_g^b)$ is compact.
\end{theorem}

This result was proven by Mumford \cite{MumfordCriteria} in the closed case. In \cite{LengthSpectraTeichMetricBoundary}, a different version of the theorem is stated. Here, we use Farb and Margalit's proof in \cite[Chap.~12]{teich} in the case $b=0$ and adapt it to the case $b>0$ to prove Theorem~\ref{theorem:mumford}. With Corollary~\ref{col:bavard}, we obtain an equivalence of Theorem~\ref{theorem:mumford} with \cite[Theorem~4.1]{LengthSpectraTeichMetricBoundary}.


\begin{proof}
    Since $\mathcal{M}(S_g^b)$ inherits the Teichmüller metric from $\mathrm{Teich}(S_g^b)$, we just need to show that $\mathcal{M}_{A, \varepsilon}(S_g^b) $ is sequentially compact. Let $(X_i)$ be a sequence in $\mathcal{M}_{A, \varepsilon}(S_g^b)$ and $\chi_i \in \mathrm{Teich}(S_g^b)$ a lift of $X_i$ for all $i$. 
    
    To show that a subsequence of $(X_i)$ converges in $\mathcal{M}_{\varepsilon}(S_g^b)$, we prove that for a fixed choice of Fenchel-Nielsen coordinates, we can choose lifts $\chi_i$ of $X_i$ to $\mathrm{Teich}(S_g^b)$ inside a rectangular compact set of the Euclidean space $\mathbb{R}_{>0}^{3g-3+2b} \times \mathbb{R}^{3g-3+b}  $. 
    
    By Theorem~\ref{theorem:Bers}, there is a pants decomposition $\mathcal{P}_i$ of $S_g^b$ such that \mbox{$\ell_{\chi_i}(\gamma) \in [\varepsilon, B]$} for all $\gamma \in \mathcal{P}_i$, where $B=\max (bA, 2\pi(2g+b-1))$.
    
    Since there are a finitely many topological types of pants decompositions of $S_g^b$, we can choose a sequence $(f_i)$ in $\mathcal{MCG}(S_g^b)$ such that, up to passing to a subsequence, $f_i(\mathcal{P}_i)=\mathcal{P}_1$. The hyperbolic structure $\Upsilon_i = f_i.\chi_i$ is also a lift of $X_i$, whose length parameters in the Fenchel-Nielsen coordinates with respect to $\mathcal{P}_1$ are between~$\varepsilon$ and~$B$.
    
    Since the Dehn twists on the curves in $\mathcal{P}_1$ change the twist parameters by $2\pi$, there is a product $h_i$ of Dehn twists on the curves of $\mathcal{P}_1$ such that the twist parameters of $h_i.\Upsilon_i$ are between $0$ and $2\pi$. Thus, the lifts $h_i.\Upsilon_i$ of $X_i$ are all inside a compact set. Therefore, there exists a converging subsequence which projects to a converging subsequence of $(X_i)$.
\end{proof}

%%%%%%

Then we state a corollary of Theorem~\ref{theorem:mumford}.

\begin{corollaire}\label{cor:mumf}
    Let $A \geqslant \varepsilon >0$. There exists a compact subset $\mathcal{Q}(A, \varepsilon) \subset \mathrm{Teich}(S_g^b)$ such that for each surface $\chi \in \mathrm{Teich}(S_g^b)$ with $\mathrm{sys}(\chi) \geqslant \varepsilon$ and boundary component lengths between $\varepsilon$ and $A$, there exists an isometric surface $\chi' \in \mathcal{Q}(A, \varepsilon)$.
\end{corollaire}

\begin{proof}
    By Theorem~\ref{theorem:mumford}, the set $\mathcal{M}_{A, \varepsilon}(S_g^b)$ of hyperbolic surfaces~$X$, up to isometry, with $\mathrm{sys}(X)\geqslant \varepsilon$ and boundary component length between $\varepsilon$ and $A$ is compact. We set $\mathcal{Q}(A, \varepsilon)$ a compact lift of $\mathcal{M}_{A, \varepsilon}(S_g^b)$ in the Teichmüller space. By definition, any surface $\chi' \in \mathrm{Teich}(S_g^b)$ with $\mathrm{sys}(\chi') \geqslant \varepsilon$ and boundary component length between $\varepsilon$ and $A$ is sent by the cover on a surface $X' \in \mathcal{M}_{A,\varepsilon}(S_g^b)$. By construction, there is a lift $\chi$ of $X'$ in $\mathcal{Q}(A, \varepsilon)$ and $\chi$ is isometric to $\chi'$.
\end{proof}


Using the analytic nature of the Teichmüller space, we also define some negligible subsets of $\mathrm{Teich}(S_g^b)$. \\

\begin{definition}
    Define the relation of equivalence $\sim$ as follows: Let $\tau$ and $\tau'$ be two arcs on $S_g^b$. We set $\tau \sim \tau'$ if for all $\chi \in \mathrm{Teich}(S_g^b)$, we have $\ell_\chi(\tau)=\ell_\chi(\tau')$. Denote $[[\tau]]$ the class of equivalence of $\tau$ for the relation~$\sim$.
\end{definition}


\begin{ex}
    By symmetry of the pair of pants, the orthogeodesics $\tau$ and $\tau'$ in Figure~\ref{fig:symExemple} are in the same class $[[\tau]]$.
    \begin{figure}[H]
        \centering
        \includegraphics[height=5cm]{figures/exempleSym.png}
        \caption{Example of two orthogeodesics in the same class of equivalence of $\sim$.}
        \label{fig:symExemple}
    \end{figure}
\end{ex}

Curves with the same property also exist and have been studied, see Chris Leininger's paper~\cite{Leininger} and references there, in part~\cite{horowitz} and~\cite{Randol}.\\

We define the length of $[[\tau]]$ as $\ell_\chi([[\tau]])=\ell_\chi(\tau)$. By definition of $\sim$, this length does not depend on the choice of representative $\tau$.

\begin{definition}
    Let $[[\tau]]$ and $[[\tau']]$ be two different equivalence classes of arcs. We define the set
    \[
        \mathcal{E}([[\tau]],[[\tau']])=\{ \chi \in \mathrm{Teich}(S_g^b) \mid \ell_\chi([[\tau]]) = \ell_\chi([[\tau']])  \}.
    \]
\end{definition}

\begin{proposition}\label{prop:E_negligible}
    The set $\mathcal{E}([[\tau]],[[\tau']])$ is a proper real analytic variety of $\mathrm{Teich}(S_g^b)$.
\end{proposition}

\begin{proof}
    By Lemma~\ref{lem:RealAnalLength}, the length function for any arc $\beta$, and thus for any class of arc $[[\beta]]$, is real analytic on $\mathrm{Teich}(S_g^b)$. Therefore, $\mathcal{E}([[\tau]],[[\tau']])$ is a real analytic variety. Moreover, by definition of $\sim$, for $[[\tau]]$ and $[[\tau']]$ different, we have 
    \[
        \mathcal{E}([[\tau]],[[\tau']]) \neq \mathrm{Teich}(S_g^b).
    \]
    
    Thus $\mathcal{E}([[\tau]],[[\tau']])$ is a proper real analytic variety of $\mathrm{Teich}(S_g^b)$.
\end{proof}

In particular, this means that $\mathcal{E}([[\tau]],[[\tau']])$ is negligible in $\mathrm{Teich}(S_g^b)$. As 
\[
    \bigcup_{[[\tau]]\neq [[\tau']]} \mathcal{E}([[\tau]],[[\tau']])
\]
is a countable union, it is also a negligible subset of $\mathrm{Teich}(S_g^b)$.






\subsection{Orthospectrum.}

Let us introduce the orthospectrum, an object analogous to the length spectrum and first defined by Basmajian in~\cite{Basmajian}.

\begin{definition}
    The \emph{orthospectrum} of a hyperbolic surface $X$ is the multiset $\mathcal{O}(X)$ of lengths of orthogeodesics on $X$, counted with multiplicities.
\end{definition}

Along with its definition, Basmajian showed the following in \cite{Basmajian}.

\begin{theorem}\label{thm:OrthoDiscrete}
    The orthospectrum is discrete.
\end{theorem}

We also define the simple orthospectrum, which is the focus of Theorem~\ref{thm:OrthoFiniteSimple}.

\begin{definition}
    The \emph{simple orthospectrum} of a hyperbolic surface $X$ is the multiset $\mathcal{O}_S(X)$ of lengths of simple orthogeodesics on $X$ counted with multiplicities.
\end{definition}



Among the properties of the orthospectrum, we highlight the following one by \cite{Basmajian}, which shows that the boundary length of a hyperbolic surface is determined by its orthospectrum. 

\begin{theorem}[Basmajian's Identity.]\label{theorem:perimetre}
    Let $X$ be a compact hyperbolic surface with geodesic boundary. Then,
    \[
        \ell (\partial X) = \sum_{\ell \in \mathcal{O}(X)} F(\ell)
    \]
    where $F(\ell)=2\sinh^{-1} \big( \frac{1}{\sinh (\ell)} \big)$. Note that $F$ is a positive decreasing function.
\end{theorem}

For the simple orthospectrum, the theorem implies that 
\[
    \ell (\partial X) > \sum_{\ell \in \mathcal{O}_S(X)} F(\ell).
\]
Thus, $X$ has at least one boundary component of length greater than $ \frac{F(osys(X))}{b}$. Moreover, it is also possible to obtain an upper bound on $\ell (\partial X)$ thanks to a corollary of \cite[Théorème~1]{Bavard}.

\begin{corollaire}\label{col:bavard}
    Let $t$ be the orthosystole of $X$. Then
    \[
        \sinh(2t)\sinh\left( \frac{\ell(\partial X)}{24g+12b-24}\right) \leqslant \frac{1}{2}.
    \]
    
\end{corollaire}
This gives us an upper bound
\[
    \ell(\partial X ) \leqslant (24g+12b-24) \sinh^{-1}\left( \frac{1}{2\sinh(2t)} \right).
\]

Another interesting result related to the simple orthospectrum is the following theorem on counting simple arcs~\cite{nickBell} inspired by Mirzakhani's Theorem on counting simple geodesics~\cite{mirzakhani}:
Two arcs are of the same type if they lie in the same $\mathcal{PMCG}(S_g^b)$-orbit.
\begin{theorem}\label{thm:comptageArc}
    Let $X$ be a hyperbolic structure with geodesic boundary on $S_g^b$. Let $\alpha_0$ be a simple essential arc on $S_g^b$. Then there exist positive constants $c(\alpha_0)$ and $\mathfrak{m}(X)$ such that
    \begin{align*}
        \lim_{L\to\infty}\frac{\# \{ \alpha \mbox{ an arc of type } \alpha_0 \mid \ell_X(\alpha) \leqslant L \}}{L^{6g-6+2b}}=c(\alpha_0)\mathfrak{m}(X).
    \end{align*}
\end{theorem}

The constant $c(\alpha_0)$ only depends on the type of the arc $\alpha_0$ and the topology of $S_g^b$. Moreover, there is a finite number of $\mathcal{PMCG}(S_g^b)$-orbits thus a finite number of classes of simple arcs~\cite{VivekaJuan}. Thus we deduce the following corollary. 

\begin{corollaire}\label{cor:comptageArc}
    Let $X$ be a hyperbolic structure with geodesic boundary on $S_g^b$. There exist positive constants $c(S_g^b)$ and $\mathfrak{m}(X)$ such that
    \begin{align}
        \lim_{L\to\infty}\frac{\# \{ t \in \mathcal{O}_S(X) \mid t \leqslant L \}}{L^{6g-6+2b}}=c(S_g^b)\mathfrak{m}(X).\label{rel:comptageArcs}
    \end{align}
\end{corollaire}

The constant $\mathfrak{m}(X)$ is the Thurston measure of the set of laminations of $S_g^b$ of length smaller than $1$. We do not need a precise definition of $\mathfrak{m}(X)$, for which we refer to~\cite{VivekaJuan}. Instead, we will use~\cite[Proposition~3.1]{mirzakhani} adapted to surfaces with boundary, which can be deduced from~\cite{MesureThurston}.

\begin{theorem}
    Let $X$ be a hyperbolic structure with geodesic boundary on $S_g^b$. Then we have
    \begin{align}
        \lim_{L\to\infty}\frac{\# \{ \gamma \in \mathcal{ML}_{g,b}(\mathbb{Z}) \mid \ell_X(\gamma) \leqslant L \}}{L^{6g-6+2b}}=\mathfrak{m}(X).\label{rel:mesureThurston}
    \end{align}
\end{theorem}

Thanks to this equality we can control the value of $\mathfrak{m}(X)$ with a lower bound on the systole of $X$ and an upper bound on the length of the boundary of $X$. Define, for the rest of this chapter, $\varepsilon_{\mathfrak{m}}>0$ to be small enough such that two non-homotopic simple closed curves shorter than $\varepsilon_{\mathfrak{m}}$ are disjoint.

\begin{proposition}\label{prop:controleMesureThurston}
    For all $g,b \geqslant 0$ and $B \geqslant 2\pi(2g+b-2)$, there exist two constants $c_1(B)$, $c_2(B)>1$ that depend only on $B$ such that for all hyperbolic structure $X$ with geodesic boundary on $S_g^b$ and with $\ell(\partial X) \leqslant B$, we have
    \begin{align*}
        \frac{1}{c_1(B)} \prod_{\gamma : \ell_X(\gamma)\leqslant \varepsilon_{\mathfrak{m}}} \frac{1}{\ell_X(\gamma)S(\ell_X(\gamma))} & \leqslant \mathfrak{m}(X) \mbox{ and,}\\
        \mathfrak{m}(X) \leqslant & c_2(B) \prod_{\gamma : \ell_X(\gamma)\leqslant \varepsilon_{\mathfrak{m}}} \left( 1 + \frac{1}{\ell_X(\gamma)S(\ell_X(\gamma))} + \frac{1}{\min\{\ell_X(\gamma),S(\ell_X(\gamma))\}} \right),
    \end{align*}
    where the products are taken over all simple closed geodesics $\gamma$ of length at most $\varepsilon_{\mathfrak{m}}$ and $S(x)=\sinh^{-1}\left( \frac{1}{\sinh(x/2)} \right)$.
\end{proposition}

This proposition is an adaptation of \cite[Proposition~3.6]{mirzakhani} to hyperbolic surfaces with boundary, whose boundary length is bounded by $B$. We prove it for the sake of completeness using the following lemma by Mirzakhani~\cite{mirzakhani}. 

\begin{lemme}
    Given $x,y,L>0$, consider the set $A_{x,y}(L)$ defined by
    \[
        A_{x,y}(L)=\{ (m,n) \in \mathbb{Z}_+ \times \mathbb{Z}_+ \mid mx + ny \leqslant L \}.
    \]
    Then for $L>6\max\{x,y\}$, we have
    \begin{align}
        \frac{1}{12}\left( \frac{L^2}{xy} \right) \leqslant \# A_{x,y}(L). \label{rel:AxyMin}
    \end{align}
    And for $L>1$, we have
    \begin{align}
         \# A_{x,y}(L) \leqslant 3L^2\left( 1 + \frac{1}{xy} +\frac{1}{\min\{x,y\}} \right). \label{rel:AxyMaj}
    \end{align}
\end{lemme}

\begin{proof}[Proof of Theorem~\ref{prop:controleMesureThurston}]
    Let $\alpha_1,...,\alpha_s$ be the curves on $X$ of length less than $\varepsilon_{\mathfrak{m}}$. By Theorem~\ref{theorem:Bers}, we can complete this family of curves into a pants decomposition $\mathcal{P}=\{ \alpha_1,...,\alpha_s,...,\alpha_{k} \}$, where $k=3g-3+b$, such that $\ell(\alpha_i) \leqslant B$. Since the Dehn twists on the curves in $\mathcal{P}$ are isometric transformations that change the twist parameters by $1$ we can assume that the twist parameters are between $0$ and $1$. Thus, by inequality~(\ref{in:combiLength}) we have
    
    \begin{multline*}
         \# \{  \gamma \in \mathcal{ML}_{g,b}(\mathbb{Z}) \mid c(B)\mathcal{L}_{\mathcal{P}}(X,\gamma) \leqslant L  \} \\
        \leqslant \# \{ \gamma \in \mathcal{ML}_{g,b}(\mathbb{Z}) \mid \ell_X(\gamma) \leqslant L \} \\ \leqslant \# \{  \gamma \in \mathcal{ML}_{g,b}(\mathbb{Z}) \mid \frac{1}{c(B)}\mathcal{L}_{\mathcal{P}}(X,\gamma) \leqslant L  \}.
    \end{multline*}
    
    Let us start by computing a lower bound for $\mathfrak{m}(X)$.\\
    
    By definition of the combinatorial length and Theorem~\ref{thm:DehnThurston}, we obtain
    {\small
    \begin{align*}
        \# \{  \gamma \in \mathcal{ML}_{g,b}(\mathbb{Z}) \mid c(B)\mathcal{L}_{\mathcal{P}}(X,\gamma) \leqslant L  \} = \# \left\{ (m_i,t_i)_{i=1}^k \in \mathcal{Z}(\mathcal{P})  \mid c(B) \left(  \sum_{i=1}^k m_i S(\ell_X(\alpha_i)) + |t_i|\ell_X(\alpha_i)  \right) \leqslant L  \right\}.
    \end{align*}
    }
    Note that we have $\{ (m_i,t_i)_{i=1}^{k} \mid t_i>0, m_i \in 2\mathbb{Z} \} \subset \mathcal{Z}(\mathcal{P})$, thus
    {\small
    \begin{align*}
        \# \left\{ (m_i,t_i)_{i=1}^k \in (\mathbb{Z}_+ \times \mathbb{Z}_+)^k  \mid c(B) \left(  \sum_{i=1}^k 2 m_i S(\ell_X(\alpha_i)) + t_i\ell_X(\alpha_i)  \right) \leqslant L  \right\}  
        & \leqslant \# \left\{ \gamma \in \mathcal{ML}_{g,b}(\mathbb{Z}) \mid \ell_X(\gamma) \leqslant L \right\} \\
        \# \left\{ (m_i,t_i)_{i=1}^k \in (\mathbb{Z}_+ \times \mathbb{Z}_+)^k \mid c(B) \left( 2 m_i S(\ell_X(\alpha_i)) + t_i\ell_X(\alpha_i)  \right) \leqslant \frac{L}{k}  \right\}  
        & \leqslant  \\
       \prod_{i=1}^{k} \# \left\{ (m_i,t_i) \in  \mathbb{Z}_+ \times \mathbb{Z}_+ \mid   2 m_i S(\ell_X(\alpha_i)) + t_i\ell_X(\alpha_i)   \leqslant \frac{L}{c(B)k}  \right\}  
        & \leqslant  
    \end{align*}
    }
    By Inequality~(\ref{rel:AxyMin}), for $L$ large enough, we have
    
    \begin{align*}
         \prod_{i=1}^{k} \frac{1}{12} \left( \frac{L}{kc(B)} \right)^2 \frac{1}{2S(\ell_X(\alpha_i))\ell_X(\alpha_i)} & 
        \leqslant \# \{ \gamma \in \mathcal{ML}_{g,b}(\mathbb{Z}) \mid \ell_X(\gamma) \leqslant L \} 
    \end{align*}  
        
    We separate the product over the curves $\alpha_i$ into a product over the ones of length at most $\varepsilon_{\mathfrak{m}}$ and a product over the other ones. 
      
    \begin{align*}
         \frac{L^{6g-6+2b}\prod_{i=s+1}^k\frac{1}{S(\varepsilon_{\mathfrak{m}})B}}{(24k^2c(B)^2 )^{3g-3+b}} \prod_{\gamma : \ell_X(\gamma)\leqslant \varepsilon_{\mathfrak{m}}}\frac{1}{\ell_X(\gamma)S(\ell_X(\gamma))} &
        \leqslant \# \{ \gamma \in \mathcal{ML}_{g,b}(\mathbb{Z}) \mid \ell_X(\gamma) \leqslant L \}  \\
         \frac{L^{6g-6+2b}}{(24k^2c(B)^2S(\varepsilon_{\mathfrak{m}})B)^{3g-3+b}} \prod_{\gamma : \ell_X(\gamma)\leqslant \varepsilon_{\mathfrak{m}}}\frac{1}{\ell_X(\gamma)S(\ell_X(\gamma))}  &
        \leqslant 
    \end{align*}  
    We obtain
    \begin{align*}  
         \frac{1}{c_1(B)} \prod_{\gamma : \ell_X(\gamma)\leqslant \varepsilon_{\mathfrak{m}}}\frac{1}{\ell_X(\gamma)S(\ell_X(\gamma))}  &
        \leqslant \frac{\# \{ \gamma \in \mathcal{ML}_{g,b}(\mathbb{Z}) \mid \ell_X(\gamma) \leqslant L \} }{L^{6g-6+2b}}
    \end{align*}
    and letting $L$ go to infinity, we get 
    \[
         \frac{1}{c_1(B)} \prod_{\gamma : \ell_X(\gamma)\leqslant \varepsilon_{\mathfrak{m}}}\frac{1}{\ell_X(\gamma)S(\ell_X(\gamma))}  
        \leqslant \mathfrak{m}(X).
    \]
    Now, let us compute the upper bound on $\mathfrak{m}(X)$. Recall that we have
    \begin{align*}
        \# \{ \gamma \in \mathcal{ML}_{g,b}(\mathbb{Z}) \mid \ell_X(\gamma) \leqslant L \} 
        \leqslant \# \{  \gamma \in \mathcal{ML}_{g,b}(\mathbb{Z}) \mid \frac{1}{c(B)}\mathcal{L}_{\mathcal{P}}(X,\gamma) \leqslant L  \}.
    \end{align*}
    By definition of the combinatorial length and Theorem~\ref{thm:DehnThurston}, we obtain
    {\small
    \begin{align*}
        \# \{ \gamma \in \mathcal{ML}_{g,b}(\mathbb{Z}) \mid \ell_X(\gamma) \leqslant L \} 
        \leqslant & \# \left\{ (m_i,t_i)_{i=1}^k \in \mathcal{Z}(\mathcal{P}) \mid \frac{1}{c(B)}\left(  \sum_{i=1}^k m_i S(\ell_X(\alpha_i)) + |t_i|\ell_X(\alpha_i)  \right) \leqslant L  \right\} \\
          \leqslant & 2^k\# \left\{ (m_i,t_i)_{i=1}^k \in (\mathbb{Z}_+ \times \mathbb{Z}_+)^k \mid \frac{1}{c(B)}\left(  \sum_{i=1}^k m_i S(\ell_X(\alpha_i)) + t_i\ell_X(\alpha_i)  \right)\leqslant L  \right\} \\
          \leqslant & 2^k\# \left\{ (m_i,t_i)_{i=1}^k \in (\mathbb{Z}_+ \times \mathbb{Z}_+)^k \mid \frac{1}{c(B)}\left(  m_i S(\ell_X(\alpha_i)) + t_i\ell_X(\alpha_i)  \right)\leqslant L  \right\} \\
           \leqslant &  \prod_{i=1}^k 2\# \left\{ (m_i,t_i) \in \mathbb{Z}_+ \times \mathbb{Z}_+ \mid  m_i S(\ell_X(\alpha_i)) + t_i\ell_X(\alpha_i)  \leqslant c(B)L  \right\}.
    \end{align*}
    }
    By Inequality~(\ref{rel:AxyMaj}), for $L$ large enough, we have
    {\small
    \begin{align*}
         \# \{ \gamma \in \mathcal{ML}_{g,b}(\mathbb{Z}) \mid \ell_X(\gamma) \leqslant L \} \leqslant &  \prod_{i=1}^k 6(Lc(B))^2\left( 1 + \frac{1}{S(\ell_X(\alpha_i))\ell_X(\alpha_i) } + \frac{1}{\min\{S(\ell_X(\alpha_i)),\ell_X(\alpha_i)\}} \right)\\
          \leqslant & L^{6g-6+2b} \prod_{i=1}^k 6c(B)^2\left(1 + \frac{1}{S(\ell_X(\alpha_i))\ell_X(\alpha_i) } + \frac{1}{\min\{S(\ell_X(\alpha_i)),\ell_X(\alpha_i)\}} \right).
    \end{align*}
    }
    Hence,
    
    \begin{align*}
          \frac{\# \{ \gamma \in \mathcal{ML}_{g,b}(\mathbb{Z}) \mid \ell_X(\gamma) \leqslant L \}}{L^{6g-6+2b}} \leqslant &  c_2(B) \prod_{\gamma : \ell_X(\gamma)\leqslant \varepsilon_{\mathfrak{m}}} \left( 1 + \frac{1}{S(\ell_X(\alpha_i))\ell_X(\alpha_i) } + \frac{1}{\min\{S(\ell_X(\alpha_i)),\ell_X(\alpha_i)\}}  \right).
    \end{align*}
    
    Letting $L$ go to infinity, we obtain 
    \[
      \mathfrak{m}(X) \leqslant c_2(B) \prod_{\gamma : \ell_X(\gamma)\leqslant \varepsilon_{\mathfrak{m}}} \left( 1 + \frac{1}{S(\ell_X(\alpha_i))\ell_X(\alpha_i) } + \frac{1}{\min\{S(\ell_X(\alpha_i)),\ell_X(\alpha_i)\}}  \right).
    \]
\end{proof}










\subsection{Hyperbolic geometry.}

Let us state several properties on geodesics and orthogeodesics using the length function. The following result is shown in \cite[Theorem~4.2.1]{buser}.

\begin{theorem}\label{thm:NonSimpleLength}
    Let $X$ be a hyperbolic surface. Then every non-simple closed geodesic on $X$ has length greater than $1$.
\end{theorem}

By doubling the surface, we deduce
\begin{corollaire}
    Let $X$ be a hyperbolic surface. Then every non-simple orthogeodesic on $X$ has length greater than $\frac{1}{2}$.
\end{corollaire}

We also have:

\begin{lemme}[Half-collar lemma]\label{lem:Half-collar}
    Let $Y$ be a pair of pants with boundary geodesics $\gamma_1, \gamma_2, \gamma_3$. The sets
    \[
        \mathcal{C}^*[\gamma_i]=\{ p \in Y \mid \sinh{(\mathrm{dist}(p,\gamma_i))}\sinh{(\tfrac{1}{2}\ell(\gamma_i))} \leqslant 1 \}
    \]
    for $i=1,2,3$ are pairwise disjoint and each of them is homeomorphic to a cylinder.
 
\end{lemme}


\begin{lemme}[Collar lemma]\label{lemme:collier}
    Let $\alpha$ and $\beta$ be two distinct closed geodesics on a hyperbolic surface $X$ such that $\alpha \cap \beta \neq \emptyset $. If $\beta$ is simple, then
    \[
        \sinh\left( \frac{\ell(\alpha)}{2} \right)\sinh\left( \frac{\ell(\beta)}{2} \right) > 1.
    \]
\end{lemme}

The proofs of Lemmas~\ref{lem:Half-collar} and~\ref{lemme:collier} can be found in~\cite{buser}. We can state a version of the collar lemma with orthogeodesics instead of closed geodesics, which can be deduced by doubling the surface.
%% 

\begin{lemme}[OrthoCollar lemma]\label{lemme:orthoCollier}
    Let $\alpha$ and $\beta$ be two distinct orthogeodesics on a hyperbolic surface $X$ such that $\alpha \cap \beta \neq \emptyset $. If $\beta$ is simple, then
    \[
        \sinh\left( \ell(\alpha) \right)\sinh\left( \ell(\beta) \right) > 1.
    \]
\end{lemme}



Now, let us state hyperbolic trigonometry formulas that we will need in the various proofs of this paper.

\begin{lemme}\label{lem:Tri+Hexa}
    For any right-angled hexagon with consecutive sides $\beta, c, \alpha, b, \gamma, a$, we have 
    \begin{align}
        \cosh(c)=\sinh(a)\sinh(b)\cosh(\gamma) - \cosh(a)\cosh(b). \label{rel:HexDroit}
    \end{align}

    For every trirectangle with sides labelled as in Figure~\ref{fig:Trirectangle}, the following relation is true:
    \[
        \cos{(\varphi)}=\tanh{(\sigma)}\tanh{(\tau)}.
    \]

    \begin{figure}[H]
        \centering
        \includegraphics[height=4cm]{figures/PentaTri2.png}
        \caption{Right-angled hexagon and trirectangle}
        \label{fig:Trirectangle}
    \end{figure}

    
\end{lemme}

The proofs can be found in \cite{buser}.\\
As already mentioned, any right-angled hexagon is determined by the lengths of three pairwise disjoint sides. A right-angled octagon can be obtained by gluing two right-angled hexagons along one side. Thus, any right-angled octagon is determined by the length of four pairwise disjoint sides and the length of one orthogonal arc between opposite sides. In the following lemma, we will see how to compute the length of the orthogonal between the last two other sides.

\begin{lemme}[Right-angled octagon]\label{lem:OctaDroit}
    We define the function $f_{\mathrm{octa}} : \mathbb{R}^5_+ \to \mathbb{R}_+$ given by
    \begin{multline*}
        f_{\mathrm{octa}}(a,x,y,z,t)  =  -\cosh(x)\cosh(t) + \sinh(x)\sinh(t) \times \\ 
        \cosh \bigg[ \cosh^{-1}\left( \frac{\cosh(y)+\cosh(x)\cosh(a)}{\sinh(x)\sinh(a)} \right) \\
        + \cosh^{-1}\left( \frac{\cosh(y)+\cosh(t)\cosh(a)}{\sinh(t)\sinh(a)} \right) \bigg]. 
    \end{multline*}
    For any right-angled octagon with four disjoint sides $\delta_1, \delta_2, \delta_3, \delta_4$ and two orthogonal arcs $\alpha$ and $\beta$ joining two opposite sides different from the $\delta_i$ as in Figure~\ref{fig:DemoOctagon}, we have:
    \[
        \cosh(\beta)  = f_{\mathrm{octa}}(\alpha,\delta_1,\delta_2,\delta_3,\delta_4).
    \]
\end{lemme}

\begin{proof}
    We can decompose the octagon into two right-angled hexagons with three disjoint sides $(\beta, \delta_4, \delta_1)$ and $(\beta, \delta_2, \delta_3)$. We can also decompose it into two other right-angled hexagons with three disjoint sides $(\alpha, \delta_1, \delta_2)$ and $(\alpha, \delta_3, \delta_4)$. The arc $\alpha$ decomposes the side between $\delta_1$ and $\delta_4$ into two segments $x$ and $x'$ such that $x$ is a side of the hexagon $(\alpha, \delta_1, \delta_2)$ and $x'$ is a side of the hexagon $(\alpha, \delta_3, \delta_4)$ as in Figure~\ref{fig:DemoOctagon}.
    \begin{figure}[H]
        \centering
        \includegraphics[height=6cm]{figures/democtagon.png}
        \caption{Right-angled octagon}
        \label{fig:DemoOctagon}
    \end{figure}
    We then apply (\ref{rel:HexDroit}) from Lemma~\ref{lem:Tri+Hexa} to the hexagons $(\alpha, \delta_1, \delta_2)$, $(\alpha, \delta_3, \delta_4)$ and $(\beta, \delta_4, \delta_1)$:
    \begin{align*}
        \cosh(x) & = \frac{\cosh(\delta_2) + \cosh(\delta_1)\cosh(\alpha)}{\sinh(\delta_1)\sinh(\alpha)}\\
        \cosh(x') & = \frac{\cosh(\delta_3) + \cosh(\delta_4)\cosh(\alpha)}{\sinh(\delta_4)\sinh(\alpha)}\\
        \cosh(\beta) & = \sinh(\delta_1)\sinh(\delta_4)\cosh(x+x') - \cosh(\delta_1)\cosh(\delta_4).
    \end{align*}
    We use the first two equalities to compute $x+x'$ and we use the expression in the last equality to conclude.
\end{proof}


In \cite[Lemma~3.2]{orthoSys} and \cite[Lemma~3.3]{orthoSys}, Masai and McShane prove the following results.

\begin{lemme}\label{lemme:In1}
    Let $P$ be a pair of pants with boundary geodesics $\alpha, \beta, \gamma$ such that \\
    $\ell(\beta) \leqslant \ell(\alpha)$. Let $\tau$ be the unique simple orthogeodesic with both endpoints on $\gamma$ as in Figure \ref{fig:dess_1}. Then, we have
    \[
        \sinh(\ell(\tau)/2) \leqslant \frac{ \cosh(\ell(\alpha)/2) }{ \sinh(\ell(\gamma)/4) }.
    \]
    Let $\tau'$ be the unique orthogeodesic with one self-intersection, both endpoints on $\gamma$ and that goes around $\alpha$ exactly once as in Figure~\ref{fig:dess_1}. Then, we have
    \begin{align}
        \ell(\tau') &  \leqslant \ell(\gamma) +2\ell(\tau), \label{In2:rel1}\\
        \mbox{and }\cosh\left( \frac{\ell(\tau')}{2}\right) & = 2 \cosh\left( \frac{\ell(\alpha)}{2}\right) \cosh\left( \frac{\ell(\tau)}{2}\right). \label{In2:rel2}
    \end{align}
    \begin{figure}[H]
        \centering
        \includegraphics[height=5cm]{figures/dessin1.png} 
        \caption{Orthogeodesics $\tau$ and $\tau'$ in $P$.}
        \label{fig:dess_1}
    \end{figure}
\end{lemme}

%%%%%%%%%





















\section{A finite number of surfaces with the same orthospectrum or simple orthospectrum}\label{sec:2}

We recall that $S_g^b$ is a genus $g$ surface of negative Euler characteristic with $b$ boundary components. We label the boundary components $\gamma_0, \gamma_1,...,\gamma_{b-1}$. Let us define, for any hyperbolic structure $X$ on $S_g^b$, the sets
\[
    \mathrm{I}(X)=\{ Y \in \mathcal{M}(S_g^b) \text{ such that } \mathcal{O}(Y)=\mathcal{O}(X) \}.
\] 
and
\[
    \mathrm{I}_S(X)=\{ Y \in \mathcal{M}(S_g^b) \text{ such that } \mathcal{O}_S(Y)=\mathcal{O}_S(X) \}.
\] 

The goal of this section is to show that each of these sets is finite.\\

In 2023, Masai and McShane showed the following theorem~\cite[Theorem~5.1]{orthoSys}:

\begin{theorem}\label{thm:OrthoFinite1}
    Let $X$ be a hyperbolic structure on $S_g^1$ with geodesic boundary. Then the set $\mathrm{I}(X)$ is finite.
\end{theorem}

Their result is restricted to surfaces with one boundary component but it can be extended to surfaces with any number of boundary components $b$.

\begin{theorem}\label{thm:OrthoFiniteb}
    Let $X$ be a hyperbolic structure on $S_g^b$ with geodesic boundary. Then, the set $\mathrm{I}(X)$ is finite.
\end{theorem}

We will only detail parts of the proof that diverge from the proof of Masai and McShane. In particular, the proof can be separated into two steps: First, show that $\mathrm{I}(X)$ is included in a compact of the form $\mathcal{M}_{A, \varepsilon}(S_g^b)$. Then show that $\mathrm{I}(X)$ is itself compact and discrete thus finite. The difference between the proof of Theorem~\ref{thm:OrthoFinite1} and Theorem~\ref{thm:OrthoFiniteb} lies in the first step, thus we will only show the first step of the proof. However, note that the second step of the proof in the case of the simple orthospectrum also works here.


\begin{proof}
   
    To prove that $\mathrm{I}(X)$ is included in a compact set, we compute two constants $A>0$ and $\varepsilon>0$ that depend on the orthospectrum of $X$ such that $\mathrm{I}(X)\subset \mathcal{M}_{A, \varepsilon}(S_g^b)$ (which is compact by Theorem~\ref{theorem:mumford}).\\
   
    If $X$ has a single boundary component $\gamma_0$ then, by Theorem~\ref{theorem:perimetre}, its length is determined by the orthospectrum of $X$ and we can choose $A=\ell(\gamma_0)=\sum_{\ell \in \mathcal{O}(X)} F(\ell)$.
    When $X$ has $b$ boundary components it is still possible to find an upper bound on their length by Corollary~\ref{col:bavard} and we can choose $A=(24g+12b-24) \sinh^{-1}\left( \frac{1}{2\sinh(2 \mathrm{osys}(X))} \right)$.\\
   
    Then Masai and McShane argue by induction to compute a lower bound on the systole of any surface $Y \in \mathrm{I}(X)$. Their argument uses an upper bound on the boundary length, the discreteness of the orthospectrum and a lower bound on the length of at least one of the boundary components. In the proof of Masai and McShane, the length of $\gamma_0$ is determined by the orthospectrum. In the case of a surface with $b$ boundary components, the length of each of the boundary components is bounded by $A$. By Theorem~\ref{theorem:perimetre}, at least one of the boundary components has a length greater than $\varepsilon_0 = \frac{\sum_{\ell \in \mathcal{O}(X)} F(\ell)}{b}$. Then, we follow the same steps as Masai and McShane to compute a lower bound $\varepsilon$ on the systole as well as a lower bound on the length of each boundary component. This is possible because each induction step does not depend on whether if the boundary components of the pair of pants correspond to non-peripheral curves on $Y$ or to boundary components of $Y$.
\end{proof}





In the case of the simple orthospectrum we can prove a result analogous to Theorem~\ref{thm:OrthoFinite1}:

\begin{theorem}\label{thm:OrthoFiniteSimple}
    Let $X$ be a hyperbolic structure on $S_g^b$. Then the set $\mathrm{I}_S(X)$ is finite.
\end{theorem}

Similarly, the proof can be done in the same two steps. However the way to prove each of these steps uses completely different techniques. Indeed, Masai and McShane's proof uses orthogeodesics with one self-intersection, which we cannot do in the case of the simple orthospectrum. Before beginning the proof, let us state a proposition that we will need for it.

\begin{proposition}\label{prop:Hummel}
    Set $A_0, \varepsilon_0 >0$, $0<k\leqslant b$ and let $(\chi_n)_{n\in\mathbb{N}} \subset \mathrm{Teich}(S_g^b)$ be a sequence of hyperbolic structures on $S_g^b$ such that for all $n$, $\mathrm{sys(\chi_n)}, \ell_{\chi_n}(\gamma_0),...,\ell_{\chi_n}(\gamma_{k}) \geqslant \varepsilon_0$ and for all $0 \leqslant j < b$ we have $\ell_{\chi_{n}}(\gamma_j) \leqslant A_0$. Moreover, for $k<j<b$, let $(\chi_n)_{n\in\mathbb{N}}$ be such that $\lim_{n \to \infty} \ell_{\chi_n}(\gamma_j)=0$. \\
    %%%
    Up to a subsequence, there exist a sequence of real number $\varepsilon_n>0$, such that the\\ $\varepsilon_n$-thick part of $\chi_n$ contains every simple orthogeodesic whose endpoints are in $\bigcup_{0 \leqslant j \leqslant k}\gamma_j$ and does not contain any $\gamma_j$ for $k < j < b$, and such that for all $\varepsilon>0$ as small as we want, we have the following.\\
    %%
    There is a subsequence $(\chi_{n_i})_{i\in\mathbb{N}}$ such that the $\varepsilon_{n_i}$-thick parts of the $\chi_{n_i}$ are\\ $(1+\varepsilon)$-bilipshitz.
\end{proposition}

This proposition results from Proposition~\cite[5.1~Chap.IV]{Hummel}.


\begin{proof}[Proof of Theorem~\ref{thm:OrthoFiniteSimple}]\label{proof:FiniSimple}

    %%%%%%%%%%%%%%%%%%%%%Compacité
    
    First, let us show that $\mathrm{I}_S(X)$ is included in $\mathcal{M}_{A, \varepsilon}(S_g^b)$ by computing for all $Y \in \mathrm{I}_S(X)$ an upper bound and a lower bound on the length of each boundary component and a lower bound on its systole. \\
    
    By Corollary~\ref{col:bavard}, we have a common upper bound $A$ on the length of the boundary components of any $Y \in \mathrm{I}_S(X)$.\\
    By Corollary~\ref{cor:comptageArc}, if $Y \in \mathrm{I}_S(X)$, we have
    \begin{align*}
        \mathfrak{m}(Y)c(S_g^b)=\lim_{L\to\infty}\frac{\# \{ t \in \mathcal{O}_S(Y) \mid t \leqslant L \}}{L^{6g+2b-6}}=\lim_{L\to\infty}\frac{\# \{ t \in \mathcal{O}_S(X) \mid t \leqslant L \}}{L^{6g+2b-6}}=\mathfrak{m}(X)c(S_g^b).
    \end{align*}
    As $\mathfrak{m}(Y)$ is constant, Proposition~\ref{prop:controleMesureThurston} implies that $\mathrm{sys}(Y)$ cannot go to $0$, i.e. there exists a constant $\varepsilon_1>0$ such that for all $Y \in \mathrm{I}_S(X)$, we have $\mathrm{sys}(Y) \geqslant \varepsilon_1$. \\

    
    Finally, let us compute a lower bound on the length of each boundary component. First, it follows from Basmajian's Identity that for all $Y \in \mathrm{I}_S(X)$ there exist a boundary component of length at least
    \[
        \varepsilon_2 := \frac{1}{b}\sum_{\ell \in \mathcal{O}_S(X)} F(\ell)>0 .
    \]
    
    
    %%%%%%%%
    
    By contradiction, let us suppose that there is a sequence of hyperbolic structures $(X_n)_{n\in\mathbb{N}} \subset \mathrm{I}_S(X)$ such that the length of some boundary components goes to zero as $n$ goes to infinity. Up to passing to a subsequence, we can assume that $\ell_{X_n}(\gamma_0) \geqslant \epsilon_2$ for all $n$. Set $0<k <b-2$ and let $\gamma_{k+1},...,\gamma_{b-1}$ be the boundary components whose length goes to $0$ as $n$ goes to infinity. We can assume that the length of each other boundary component has $\varepsilon_2$ as a lower bound.
    
    For each $X_n$, define $\varepsilon_n$ as in Proposition~\ref{prop:Hummel}. By this proposition, for all $\varepsilon'>0$, up to a subsequence, the $thick_{\varepsilon_n}(X_n)$ are all \\
    $(1+\varepsilon')$-bilipschitz. Fix $\varepsilon'$; then set
    \[
        N(X_n,L)= \# \{ \tau \mbox{ simple orthogeodesic in } X_n \mid \ell_{X_n}(\tau) < L \}
    \]
    and
    \[
        M(X_n,L) = \# \{ \tau \mbox{ simple orthogeodesic in } thick_{\varepsilon_n}(X_n) \mid \ell_{X_n}(\tau) < L \}.
    \]
    Note that, as we already have a lower bound on the length of $\gamma_0$, there is at least one boundary component whose length does not go to $0$ as $n$ tends to infinity. Thus for $L$ sufficiently large, $M(X_n,L)$ is not empty when $N(X_n,L)$ is not empty.\\
    
    By Theorem~\ref{thm:comptageArc}, there exists $L_{\varepsilon'}$ large enough such that
    \begin{align}
        N \left( X_0,\frac{L_{\varepsilon'}}{(1+\varepsilon')} \right) > & (1-\varepsilon')A'\mathfrak{m}(X_0)\left( \frac{L_{\varepsilon'}}{(1+\varepsilon')} \right)^{6g-6+2b} \label{rel:comptage1} \\
       \mbox{and } M \left( X_0,L_{\varepsilon'} \right) < & (1+\varepsilon')B'\mathfrak{m}(X_0) L_{\varepsilon'} ^{6g-6+2b}, \label{rel:comptage2}
    \end{align}
    where $A'$ is the sum of $c(\alpha_0)$ where $\alpha_0$ varies among all types of arcs in $X_0$ and $B'$ is the sum of $c(\alpha_0)$ where $\alpha_0$ varies among the types of arcs in $X_0$ with their endpoints in $\bigcup_{0 \leqslant j < k}\gamma_j$.\\
    
    For $n$ large enough, any arc with at least one endpoint on a boundary component among $\gamma_{k+1}$,...,$\gamma_{b-1}$ has length greater than $\frac{L_{\varepsilon'}}{1+\varepsilon'} $. In part we have 
    \[
        N \left( X_n, \frac{L_{\varepsilon'}}{(1+\varepsilon')} \right) = M\left( X_n,\frac{L_{\varepsilon'}}{(1+\varepsilon')} \right).
    \]
    As $\mathcal{O}_S(X_n)=\mathcal{O}_S(X)$ for all $n$, we have $N \left( X_n, \frac{L_{\varepsilon'}}{(1+\varepsilon')} \right) = N \left( X_0, \frac{L_{\varepsilon'}}{(1+\varepsilon')} \right)$. Moreover, for all $n$ we have $thick_{\varepsilon_n}(X_n)$ and $thick_{\varepsilon_0}(X_0)$ are $(1+\varepsilon')$-bilipschitz, thus $M\left( X_n,\frac{L_{\varepsilon'}}{(1+\varepsilon')} \right) \leqslant M(X_0,L_{\varepsilon'})$. We obtain
    \[
        N \left( X_0, \frac{L_{\varepsilon'}}{(1+\varepsilon')} \right) \leqslant M(X_0,L_{\varepsilon'}).
    \]
    By (\ref{rel:comptage1}) and (\ref{rel:comptage2}), we have
    \[
        (1-\varepsilon')A'\mathfrak{m}(X_0)\left( \frac{L_{\varepsilon'}}{(1+\varepsilon')} \right)^{6g-6+2b} < (1+\varepsilon')B'\mathfrak{m}(X_0) L_{\varepsilon'}^{6g-6+2b}.
    \]
    Thus, for all $\varepsilon'>0$, we have
    \[
        A' (1-\varepsilon') < B'(1+\varepsilon')^{5-6g-2b}.
    \]
    Letting $\varepsilon'$ tend to $0$ we obtain that $A' \leqslant B'$, a contradiction. Thus, there exists a lower bound $\varepsilon_3$ on the length of each boundary component.\\
    If we set $\varepsilon = \min \{ \varepsilon_1, \varepsilon_2, \varepsilon_3 \}$, then $\mathrm{I}_S(X) \subset \mathcal{M}_{A, \varepsilon}(S_g^b)$, which is compact.\\
    
    %%%%%%%%%%%%%%%%%%%%%%%%%%%%%%
    
    For the second step of the proof, let us fix a hexagon decomposition $\mathcal{H}$ on $S_g^b$ and set $\varphi_{\mathcal{H}}$ as in Theorem~\ref{HexaCoord}.

    Because $\mathrm{I}_S(X)$ is included in a compact subset of $\mathcal{M}(S_g^b)$, there is a lift $\Tilde{\mathrm{I}}_S(X)$ of $\mathrm{I}_S(X)$ in $\mathrm{Teich}(S_g^b)$ which is also included in a compact~$\mathrm{C}$. We will show that $\Tilde{\mathrm{I}}_S(X)$ is finite, and so is $\mathrm{I}_S(X)$.
    
    Since $\Tilde{\mathrm{I}}_S(X)$ is included in a compact set, for any sequence $\chi_n \in \Tilde{\mathrm{I}}_S(X)$, there is a subsequence such that $\chi_n \to \chi \in \mathrm{C}$. The map $\varphi_{\mathcal{H}}$ is continuous, so \mbox{$ \varphi_{\mathcal{H}}(\chi_n) \to \varphi_{\mathcal{H}}(\chi)$}. We recall that for every $n$, we have $\mathcal{O}_S(\chi_n)=\mathcal{O}_S(X)$ and that by Theorem~\ref{thm:OrthoDiscrete} the orthospectrum is discrete. 
    
    As, $\varphi_{\mathcal{H}}(\chi_n) \subset \mathcal{O}_S(X)$, we deduce from Wolpert's lemma \ref{lemme:OrthoWolpert} that there exists $N \in \mathbb{N}^*$ such that for all $n', m' >N$, $\varphi_{\mathcal{H}}(\chi_{n'})=\varphi_{\mathcal{H}}(\chi_{m'})$. Thus, for all $n', m' > N$, $\chi_{n'} = \chi_{m'} = \chi $ and $\chi \in \Tilde{\mathrm{I}}_S(X)$. So $\Tilde{\mathrm{I}}_S(X)$ is finite, and so is $\mathrm{I}_S(X)$. 
    
\end{proof}




















\section{Generic surfaces are determined by their orthospectrum and by their simple orthospectrum}\label{sec:3}

In this section, we are going further in our characterization of the rigidity of the orthospectrum and the simple orthospectrum. We prove the following theorem:

\begin{theorem}\label{thm:Main2}
    Let $\mathcal{V}_g^b$ be the subset of all $\chi \in \mathrm{Teich}(S_g^b)$ for which there exists\\
    $\Upsilon \in \mathrm{Teich}(S_g^b)$ non-isometric to $\chi$ such that $\mathcal{O}(\chi)=\mathcal{O}(\Upsilon)$. Similarly, let $\mathcal{W}_g^b$ be the subset of all $\chi \in \mathrm{Teich}(S_g^b)$ for which there exists $\Upsilon \in \mathrm{Teich}(S_g^b)$ non-isometric to $\chi$ such that $\mathcal{O}_S(\chi)=\mathcal{O}_S(\Upsilon)$. 
    
    Then, both $\mathcal{V}_g^b$ and $\mathcal{W}_g^b$ are negligible sets of $\mathrm{Teich}(S_g^b)$. Moreover, $\mathcal{W}_g^b$ is a proper local real analytic subvarieties of $\mathrm{Teich}(S_g^b)$.
\end{theorem}

This result is a version of Wolpert's Theorem \cite{Wolpert} for the orthospectrum and the simple orthospectrum instead of the length spectrum. We adapt the proof given by Buser in \cite[Chap. 10]{buser} to our case. Doing so, we do not require non-simple arcs, thus most of the proof is the same both for the orthospectrum and the simple orthospectrum. However there are still some key differences that will be highlighted.\\


We first set up all the context and notation, then we prove a lemma on the first lengths of the orthospectrum and simple orthospectrum, and finally we prove Theorem~\ref{thm:Main2}.













\subsection{Set up and prerequisites}\label{subsec:SetUp}

First let us state two propositions on two different kinds of open subsets of the Teichmüller space that will play a key role in the proof of Theorem~\ref{thm:Main2}. Each kind corresponds to the orthospectrum case and the simple orthospectrum case. These subsets will play the same role in the proof of Theorem~\ref{thm:Main2} but the subtle differences between the two will influence the end of the proof.\\


In the case of the orthospectrum, there exist open balls in $\mathrm{Teich}(S_g^b)$ with the following property.

\begin{proposition}\label{prop:Ux}
    For all $\chi \in \mathrm{Teich}(S_g^b) \setminus \bigcup_{[[\tau]]\neq [[\tau']]} \mathcal{E}([[\tau]],[[\tau']])$, there exists a constant $q_\chi>1$ such that the open ball $\mathcal{U}_\chi$ of center $\chi$ and radius $\frac{\log(q_\chi)}{2}$ satisfies the following property: there exist constants $\varepsilon_\chi>0$ and $A_\chi>0$ such that for all $\Upsilon \in \mathcal{U}_\chi$ and $\Upsilon' \in \mathrm{Teich}(S_g^b)$ with $\mathcal{O}(\Upsilon')=\mathcal{O}(\Upsilon)$, there is a surface isometric to $\Upsilon'$ in $\mathcal{Q}(A_\chi,\varepsilon_\chi)$.
\end{proposition}



\begin{proof}
    First, let us define the following:
    \[
        \delta(L)=\min_{t,t'\in\mathcal{O}(\chi)}\left\{t'-t \mid t<t' \leqslant L\right\}.
    \]
    This quantity measures the smallest gap between lengths of orthogeodesics of length bounded by $L$. Since there are finitely many orthogeodesics of length shorter than any $L$, $\delta(L)$ is well defined for any $L>0$ greater than the second shortest length in $\mathcal{O}(\chi)$. 
    
    To deploy this proof we first need to recall the details of the induction parts of the proof of Theorem~\ref{thm:OrthoFiniteb}, especially, how to compute a lower bound on the systole and the length of the boundary components of a surface $\chi'$ with its orthospectrum. Let $\chi' \in \mathrm{Teich}(S_g^b)$ such that $\mathcal{O}(\chi')=\mathcal{O}(\chi)$. \\
    
    By Corollary~\ref{col:bavard}, if $A=(24g+12b-24) \sinh^{-1}\left( \frac{1}{2\sinh(2 \mathrm{osys}(\chi))} \right)$, then $A \geqslant \ell(\partial \chi')$. By Theorem~\ref{theorem:perimetre}, at least one of the boundary components of $\chi'$ has a length greater than $\varepsilon_0 = \frac{\sum_{\ell \in \mathcal{O}(\chi)} F(\ell)}{b}$. We choose this boundary component to be $\gamma_0$.\\
   
   Define $\mathcal{P}$ a pants decomposition on $\chi'$ as in Theorem~\ref{theorem:Bers}. The curves in $\mathcal{P}$ and the boundary components of $\chi'$ have length at most $B=\max\{ A, \mathrm{area}(\chi)\}$. \\
   
   Construct a rooted graph $G$ dual to the pants decomposition, as follows. Each vertex corresponds to a pair of pants in $\chi'$ given by the pants decomposition $\mathcal{P}$. The root corresponds to the pair of pants, denoted by $P_0$, with $\gamma_0$ as one of its boundary components. A pair of vertices is joined by an edge each time the corresponding pairs of pants have a boundary component in common.
   
    \begin{figure}[H]
        \centering
        \includegraphics[height=6cm]{figures/Graphe.png}
        \caption{Example of a rooted graph $G$.}
        \label{fig:Graphe}
    \end{figure}
    
    Choose a spanning rooted tree $T$ in $G$. We show by induction on the distance to the root that we have a lower bound on the length of each boundary component of each pair of pants corresponding to the vertices.
    
    The base case starts with $P_0$. Label $\alpha_1$ and $\beta_1$ the other two boundary components of $P_0$. With Lemma~\ref{lemme:In1} we compute an upper bound $L_0$ on two specific orthogeodesics $\tau_0$ and $\tau_0'$ (see Figure~\ref{fig:dess_9bis}) that depends on $\varepsilon_0$ and $B$. 
    \begin{align}
        \ell(\tau_0) & \leqslant 2\sinh^{-1} \left(  \frac{\cosh(B/2)}{\sinh(\ell(\varepsilon_0))/4} \right)\\ 
        \ell(\tau_0') & \leqslant 4\sinh^{-1} \left(  \frac{\cosh(B/2)}{\sinh(\ell(\varepsilon_0))/4} \right)+B=L_0.\label{rel:L0}
    \end{align}
    
    \begin{figure}[H]
        \centering
        \includegraphics[height=5cm]{figures/dessin9bis.png}
        \caption{The pair of pants $P_0$ and the two orthogeodesics used to compute a lower bound on $\ell(\alpha_1)$.}
        \label{fig:dess_9bis}
    \end{figure}
    
    By (\ref{In2:rel1}), we have
    \begin{align*}
        \ell(\alpha_1 ) = 2 \cosh^{-1}\left(    \frac{ \cosh( \ell(\tau_0')/2 )  }{ \cosh( \ell(\tau_0)/2 ) }     \right) .
    \end{align*}
    
    Remark that the function $f(t,d)=2\cosh^{-1} \left(  \frac{\cosh((t+d)/2)}{\cosh(t/2)}  \right)$ is increasing in $t$ and $d$, thus 
    \begin{align}
         2\cosh^{-1} \left(  \frac{\cosh((\mathrm{osys}(\chi)+ \delta' )/2)}{\cosh(\mathrm{osys}(\chi)/2)}  \right) \leqslant   2\cosh^{-1} \left(  \frac{\cosh(\ell_\chi(\tau)/2+ (\ell_\chi(\tau')-\ell_\chi(\tau)) /2)}{\cosh(\ell_\chi(\tau)/2)}  \right)  \label{in:preuveUx}
    \end{align}
    for all orthogeodesics $\tau$ and $\tau'$ in $\mathcal{O}(\chi)$ and $\delta'$ such that $\ell_\chi(\tau) < \ell_\chi(\tau') \leqslant L$ and 
    \[
        \delta' \leqslant \min\{ \ell_\chi(\tau') - \ell_\chi(\tau) \mid \ell_\chi(\tau) < \ell_\chi(\tau') \leqslant L\}.
    \]
     Thus we have a lower bound $\varepsilon_1$ on the length of $\alpha_1$ as follows:
     \[
        \varepsilon_1 = 2\cosh^{-1} \left( \frac{   \cosh( (\mathrm{osys}(\chi) + \delta(L_{0}) )/2 ) }{  \cosh( \mathrm{osys}(\chi)/(2))}  \right) \leqslant \ell(\alpha_1).
    \]
    By symmetry, $\varepsilon_1$ is also a lower bound on the length of $\beta_1$.\\
    
    
    
    At step $i>0$, we look at any vertex in $T$ at distance $i$ from the root. Take a path from the root to this vertex, we consider the union of pair of pants that corresponds to vertices on this path, see Figure~\ref{fig:corail}. Then we compute an upper bound $L_i$ on the length of two orthogeodesics $\tau_i$ and $\tau_i'$ as in Figure~\ref{fig:corail}.
    \begin{align}
        \ell(\tau_i) & \leqslant 2iB + 2\sinh^{-1} \left(  \frac{\cosh(B/2)}{\sinh(\varepsilon_i)/4} \right) + \sum_{j=1}^i \cosh^{-1} \left( \frac{\cosh(B/2)(1+\cosh(B/2))}{\sinh(\varepsilon_{j-1}/2)\sinh(\varepsilon_{j}/2)} \right)\\ 
        \ell(\tau_i') & \leqslant (2i+1)B + 4\sinh^{-1} \left(  \frac{\cosh(B/2)}{\sinh(\varepsilon_i)/4} \right) + \sum_{j=1}^i \cosh^{-1} \left( \frac{\cosh(B/2)(1+\cosh(B/2))}{\sinh(\varepsilon_{j-1}/2)\sinh(\varepsilon_{j}/2)} \right)=L_i. \label{rel:Li}
    \end{align}
    
    \begin{figure}[H]
        \centering
        \includegraphics[height=6cm]{figures/Corail.png}
        \caption{Subsurface of $Y$ made out of the different pairs of pants and the orthogeodesics $\tau_i$ and $\tau_{i}'$.}
        \label{fig:corail}
    \end{figure}
    
    By (\ref{In2:rel1}), we have
    \begin{align*}
        \ell(\alpha_{i+1} ) = 2 \cosh^{-1}\left(\frac{ \cosh( \ell(\tau_i')/2 )  }{  \cosh( \ell(\tau_i)/2 ) } \right) .
    \end{align*}
    
    Then, by Inequality~\ref{in:preuveUx}
    \[ 
        \varepsilon_{i+1} = 2\cosh^{-1} \left( \frac{   \cosh( (\mathrm{osys}(\chi) + \delta(L_{i}) )/2 ) }{  \cosh( \mathrm{osys}(\chi)/(2))}  \right) \leqslant \ell(\alpha_{i+1}).
    \]
    By symmetry, $\varepsilon_{i+1}$ is also a lower bound on the length of $\beta_{i+1}$.\\
    
    There is a finite number of vertices in $G$, so there is only a finite number of steps. By computing a lower bound on the length of the boundary components of each pair of pants corresponding to a vertex in $G$, we obtain a lower bound on the systole and the length of each boundary components of $\chi'$.\\
    
    Let $n$ be diameters of $T$ and $L_{0}(\chi), L_{1}(\chi),...,L_{n}(\chi)=L(\chi)$ the different upper bounds computed along the induction process for $\chi$. \\
    
    In the next part we construct $q_\chi$ by following each step of the induction process of the previous part of the proof. At each step of this process we also compute upper and lower bounds which will help control the various upper and lower bounds obtained during the induction process for $\Upsilon'$.\\
    

    Denote $\mathcal{O}(\chi)\cap [0,L(\chi)]=\{a_1,...,a_k\}$ and $a_{k+1}=\min\{ a \in \mathcal{O}(\chi) \mid a > a_k\}$. 
    Then define 
    \[
        \delta_i= \min_{1 \leqslant j \leqslant k-1}\{a_{j+1}-a_j \mid a_{j+1} \neq a_j \mbox{ and } a_j \leqslant L_i(\chi)\}. 
    \]
    Since $\chi \in \mathrm{Teich}(S_g^b) \setminus \bigcup_{[[\tau]]\neq [[\tau']]} \mathcal{E}([[\tau]],[[\tau']])$, each length that appears several times in $\mathcal{O}(\chi)$ is realized by an arc lying in the same equivalence class of~$\sim$.  \\
    
    Set $A_q = (24g+12b-24) \sinh^{-1}\left( \frac{1}{2\sinh(2 \mathrm{osys}(\chi)/q)} \right)$, $B_q=\max\{ A_q, \mathrm{area}(\chi)\}$ and
    \[
        \delta_q(L)=\min_{1 \leqslant j \leqslant k-1}\left\{\frac{a_{j+1}}{q}-qa_j \mid a_{j+1}\neq a_j \mbox{ and } a_j \leqslant L\right\}.
    \]
    
    
    We follow the induction process of the previous part of the proof. Let us start with $i=0$.\\
    Set $\varepsilon_{q_0}=F(q_0\cdot \mathrm{osys}(\chi))/b$ and
    \[
        L_{q_0}= 4\sinh^{-1} \left( \frac{\cosh(B_{q_0}/2)}{\sinh(\varepsilon_{q_0})/4} \right) +B_{q_0}
    \]
    Set $a_{j_0}=\min\{ a_j \mid a_j > L_{0}(\chi) \}$ and let $q_0>1$ be small enough such that 
    \[ 
         L_{q_0}<\frac{a_{j_0}}{q_0}.
    \]
    Such a constant $q_0$ does exist: indeed, as $q_0$ goes to $1$, we have $L_{q_0} \to L_0(\chi)$, $\frac{a_{j_0}}{q_0} \to a_{j_0}$ and $L_0(\chi) < a_{j_0}$.\\
    Then at $i>0$, set
    \[
        \varepsilon_{q_{i}} = 2\cosh^{-1} \left( \frac{   \cosh( (\mathrm{osys}(\chi)/q_{i} + \delta_{q_{i}}(L_{q_{i-1}}) )/2 ) }{  \cosh( \mathrm{osys}(\chi)/(2q_{i}))}  \right)
    \]
    and
    \[
        L_{q_{i}} = 2i\cdot B_{q_{i}} + 2\sinh^{-1}\left( \frac{\cosh(B_{q_{i}}/2)}{\sinh(\varepsilon_{q_{i}})/4} \right) + \sum_{l=1}^{i} \cosh^{-1} \left(  \frac{ \cosh(B_{q_{i}}/2)(1+\cosh(B_{q_{i}}/2))  }{ \sinh(\varepsilon_{q_{l-1}})\sinh(\varepsilon_{q_l}) } \right).
    \]
    set $a_{j_{i}}= \min( a_j \mid a_j > L_{i}(\chi) ) $ and let $q_{i}>1$ small enough such that 
    \[ 
         L_{q_{i}}<\frac{a_{j_{i}}}{q_{i}}. 
    \]
    Such a constant $q_{i}$ does exist: indeed, as $q_{i}$ goes to $1$ we have $L_{q_{i}} \to L_{i}(\chi)$, $\frac{a_{j_{i}}}{q_{i}} \to a_{j_{i}}$ and $L_{i}(\chi) < a_{j_{i}}$.\\
    At step $i=n$, we have that $L_{q_n}<\frac{a_{k+1}}{q_n}$.\\
    
    We set 
    \[
        q_\chi=\min \left\{q_0,...,q_n,\max_{1\leqslant j \leqslant k}\left\{q>1 \mid a_{j+1}-\frac{a_{j+1}}{q}< \frac{\delta_n}{2} \mbox{ and } qa_j - a_j < \frac{\delta_n}{2}  \right\} \right\}.
    \] 
    
    Finally, in the last part, let us prove that Proposition~\ref{prop:Ux} is true for this choice of $q_\chi$. Once again, we follow the steps of the induction process of the first part. For all $\Upsilon \in \mathcal{U}_\chi$ and for all arc or curve $\beta$ on $S_g^b$, by Wolpert's Lemma~\ref{lemme:Wolpert},~\ref{lemme:OrthoWolpert} we have
    \[
        \frac{\ell_\chi(\beta)}{q_\chi} \leqslant \ell_\Upsilon(\beta) \leqslant q_\chi \ell_\chi(\beta).    
    \]
    Set $\{a'_1,...,a'_k\}=\mathcal{O}(\Upsilon)\cap [0,\frac{a_{k+1}}{q_\chi}]=\mathcal{O}(\Upsilon')\cap [0,\frac{a_{k+1}}{q_\chi}]$ and $a'_{k+1}=\min\{ a' \in \mathcal{O}(\Upsilon) \mid a' > a'_k \}$. 
    We have $\frac{a_j}{q_\chi} \leqslant a'_j \leqslant q_\chi a_j < \frac{a_{j+1}}{q_\chi}$ for all $1 \leqslant j \leqslant k$ and $\frac{a_{k+1}}{q_\chi} \leqslant a'_{k+1}$. Moreover, as $\chi \in \mathrm{Teich}(S_g^b) \setminus \bigcup_{[[\tau]]\neq [[\tau']]} \mathcal{E}([[\tau]],[[\tau']])$ each length that appears several times in $\mathcal{O}(\chi)$ is realized each time by an arc of the same class of~$\sim$. Thus we have the same for the lengths in $\mathcal{O}(\Upsilon)$ and when $a_j=a_{j+1}$ we also have $a'_j=a'_{j+1}$.\\
    
    Denote $L_{0}(\Upsilon'), L_{1}(\Upsilon'),...,L_{n}(\Upsilon')$ the different upper bounds computed as in~(\ref{rel:L0}) and~(\ref{rel:Li}) during the induction process for $\Upsilon'$, and $\varepsilon_{0}(\Upsilon'),...,\varepsilon_{n}(\Upsilon')$ the different lower bounds. Since $\Upsilon$ and $\Upsilon'$ have the same orthospectrum, they share the same bounds. \\
    We have an upper bound $A_{\Upsilon'} = (24g+12b-24) \sinh^{-1}\left( \frac{1}{2\sinh(2 \mathrm{osys}(\Upsilon'))} \right)$ on the length of the boundary component of $\Upsilon'$ that is smaller than $A_{q_\chi}$, thus an upper bound \\
    $B_{\Upsilon'}=\max\{ A_{\Upsilon'}, \mathrm{area}(\Upsilon')\}$ smaller than $B_{q_\chi}$. Moreover we have a lower bound
    \[
        \varepsilon_{0}(\Upsilon')=F(\mathrm{osys}(\Upsilon'))/b=F(\mathrm{osys}(\Upsilon))/b
    \]
    bigger than $\varepsilon_{q_0}$ thus $L_{0}(\Upsilon') < L_{q_0} < \frac{a_{j_0}}{q_0}$.\\
    
    This implies that
    \begin{multline*}
        \delta_{\Upsilon'}^1= \min\left\{\ell_{\Upsilon'}(\tau') - \ell_{\Upsilon'}(\tau) \mid \ell_{\Upsilon'}(\tau) < \ell_{\Upsilon'}(\tau') \leqslant L_0(\Upsilon') \right\} \\
        = \min_{1 \leqslant j \leqslant k-1}\left\{a_{j+1}'-a_j' \mid a_{j+1}\neq a_j \mbox{ and } a_j \leqslant L_0(\chi)\right\}.
    \end{multline*}

    Thus, $\delta_{\Upsilon'}^1  \geqslant \delta_{q_{1}}(L_{q_0}) \geqslant \delta_{q_\chi}(L_0(\chi))$. Combined with $ \frac{\mathrm{osys}(\chi)}{q_\chi} \leqslant  \mathrm{osys}(\Upsilon)=\mathrm{osys}(\Upsilon')$, we obtain that $\varepsilon_{1}(\Upsilon') \geqslant \varepsilon_{q_{1}}$.\\
    
    Then at each step $i$, we have $\varepsilon_{i-1}(\Upsilon') \geqslant \varepsilon_{q_{i-1}}$, thus $ L_{i-1}(\Upsilon') \leqslant L_{q_{i-1}} < \frac{a_{j_{i-1}}}{q_{i-1}}$. This implies that
    \begin{multline*}
        \delta_{\Upsilon'}^{i}= \min\left\{\ell_{\Upsilon'}(\tau') - \ell_{\Upsilon'}(\tau) \mid \ell_{\Upsilon'}(\tau) < \ell_{\Upsilon'}(\tau') \leqslant L_{i-1}(\Upsilon') \right\} \\ 
        = \min_{1 \leqslant j \leqslant k-1}\left\{a_{j+1}'-a_j' \mid a_{j+1}\neq a_j \mbox{ and } a_j \leqslant L_{i-1}(\chi)\right\}.
    \end{multline*}
    
    Thus $\delta_{\Upsilon'}^{i} \geqslant \delta_{q_{i}}(L_{q_{i-1}}) \geqslant \delta_{q_\chi}(L_{i-1}(\chi))$ and  $\varepsilon_{i}(\Upsilon') \geqslant \varepsilon_{q_{i}}$.\\
    
    In conclusion, for all $\Upsilon \in \mathcal{U}_\chi$ and $\Upsilon' \in \mathrm{Teich}(S_g^b)$ such that $\mathcal{O}(\Upsilon')=\mathcal{O}(\Upsilon)$ we have
    \[
        \mathrm{sys}(\Upsilon') \geqslant \varepsilon_\chi = 2\cosh^{-1} \left( \frac{   \cosh(\mathrm{osys}(\chi)/q_\chi + \delta_{q_\chi}(L(\chi))/2 )}{  \cosh( \mathrm{osys}(\chi)/(2q_\chi) )} \right).
    \]
    Moreover, this lower bound on the systole of $\Upsilon'$ is also a lower bound on the length of each boundary component of $\Upsilon'$. We also have an upper bound $A_\chi=(24g+12b-24) \sinh^{-1}\left( \frac{1}{2\sinh(2 \mathrm{osys}(\chi)/q_\chi)} \right)$ on the length of each boundary component. By Corollary~\ref{cor:mumf}, there exists a surface isometric to $\Upsilon'$ in $\mathcal{Q}(A_\chi,\varepsilon_\chi)$.
\end{proof}


In the case of the simple orthospectrum, we have a similar property, but this time balls can be defined around any point of $\mathrm{Teich}(S_g^b)$.

%%%%%%%%%%%%%%%%%%%%%%%%%%%%%%%%%%%%%%%%%%%%%%%%%%%%
%%%%%%%%Reformuler la proposition
\begin{proposition}\label{prop:UxS}
    For all $\chi \in \mathrm{Teich}(S_g^b)$, there exists $q_\chi^S >1$ such that for $\mathcal{U}_\chi^S$ an open ball of center $\chi$ and radius $\frac{\log(q_\chi^S)}{2}$ there exist $A_\chi^S, \varepsilon_\chi^S >0$ such that for all $\Upsilon \in \mathcal{U}_\chi^S$ and $\Upsilon' \in \mathrm{Teich}(S_g^b)$ with $\mathcal{O}_S(\Upsilon')=\mathcal{O}_S(\Upsilon)$ there is a surface isometric to $\Upsilon'$ in $\mathcal{Q}(A_\chi^S,\varepsilon_\chi^S)$.
\end{proposition}



\begin{proof}
    
    For surfaces with a single boundary components, we will see along the proof that any $q_\chi^S >1$ works. For surfaces with at least two boundary components, let $A'$ be the sum of $c(\alpha_0)$ where $\alpha_0$ varies among on all types of arcs in $\chi$, with $c(\alpha_0)$ the constant as in Theorem~\ref{thm:comptageArc}. Let $B'$ be the sum of $c(\alpha_0)$ where $\alpha_0$ varies among the types of arcs in $\chi$ with no endpoints on a fixed boundary component, the component that we choose to fix does not matter as $B'$ is a topological constant. \\
    Then, by definition $A' > B'$ and we can set $q_\chi^S>1$ such that $A' \left ( \frac{1}{q_\chi^S} \right)^{6g+2b-6} > B'$.\\
    
    In both cases, define
    \begin{align*}
        A_\chi^S = & (24g+12b-24)\sinh^{-1}\left( \frac{1}{2\sinh(2\mathrm{osys}(\chi)/q_\chi^S)} \right) \\
        \varepsilon_\chi^1 = & F(q_\chi^S\mathrm{osys}(\chi)).
    \end{align*}
    
    By Wolpert's Lemma~\ref{lemme:OrthoWolpert}, for all $\Upsilon \in \mathcal{U}_\chi^S$ and $\Upsilon' \in \mathrm{Teich}(S_g^b)$ such that $\mathcal{O}_S(\Upsilon')=\mathcal{O}_S(\Upsilon)$, we have $\frac{\mathrm{osys}(\chi)}{q_\chi^S} \leqslant \mathrm{osys}(\Upsilon) = \mathrm{osys}(\Upsilon') \leqslant q_\chi^S\mathrm{osys}(\chi)$. Thus, by definition of $A_\chi^S$ and by Corollary~\ref{col:bavard} we have
    \begin{align*}
        \ell_{\Upsilon}(\gamma_j), \ell_{\Upsilon'}(\gamma_j) \leqslant A_\chi^S,
    \end{align*}
    for all $0 \leqslant j < b$.
    
    
    Moreover, by Theorem~\ref{theorem:perimetre} $\chi$ has at least one boundary component, that we denote $\gamma_0$, of length greater than $F(\mathrm{osys}(\chi))$. By Wolpert's Lemma~\ref{lemme:OrthoWolpert} we have $\ell_{\Upsilon}(\gamma_0) \geqslant \varepsilon_\chi^1$ and there is at least one boundary component of $\Upsilon'$ of length greater than $\varepsilon_\chi^1$ that we denote $\gamma_{j_0}$. Note that, even if $\Upsilon$ and $\Upsilon'$ have the same orthosystole, we cannot be sure that we can use it to get a lower bound on the length of the same boundary component, hence why we denote it $\gamma_{j_0}$ and note $\gamma_0$ in the case of $\Upsilon'$. \\
    
    Moreover $\mathrm{sys}(\Upsilon) \geqslant \frac{\mathrm{sys}(\chi)}{q_\chi^S}$ for all $\Upsilon \in \mathcal{U}_\chi^S$, so by Proposition~\ref{prop:controleMesureThurston}, 
    
    \begin{align*}
        \mathfrak{m}(\Upsilon) & \leqslant c_2(B_\chi^S) \prod_{\gamma : \ell_\Upsilon(\gamma)\leqslant \varepsilon_{\mathfrak{m}}} \left( 1 + \frac{1}{\ell_\Upsilon(\gamma)S(\ell_\Upsilon(\gamma))} + \frac{1}{\min\{\ell_\Upsilon(\gamma),S(\ell_\Upsilon(\gamma))\}} \right) \\
        & \leqslant c_2(B_\chi^S)  \left( 1 + \frac{1}{\left(\mathrm{sys}(\chi)/q_\chi^S\right) S(\varepsilon_{\mathfrak{m}})} + \frac{1}{\min\{\left(\mathrm{sys}(\chi)/q_\chi^S\right),S(\varepsilon_{\mathfrak{m}})\}} \right)^{3g+b-3}
    \end{align*}
    
    where $B_\chi^S=\max\{A_\chi^S, 2\pi(2g+b-2)\}$. For all $\Upsilon' \in \mathrm{Teich}(S_g^b)$ such that \\
    $\mathcal{O}_S(\Upsilon')=\mathcal{O}_S(\Upsilon)$, we have $\mathfrak{m}(\Upsilon')=\mathfrak{m}(\Upsilon)$, so there is an upper bound on $\mathfrak{m}(\Upsilon')$ that only depends on $\chi$ and $q_\chi^S$. If $\mathrm{sys}(\Upsilon')$ goes to $0$ then $\mathfrak{m}(\Upsilon')$ goes to infinity, which is impossible. So there exists a lower bound $\varepsilon_\chi^2$ on $\mathrm{sys}(\Upsilon')$.\\ 
    
    
    If $b>1$ we still need to compute a lower bound on the length of the boundary components of $\Upsilon'$ different from $\gamma_{j_0}$. We proceed by contradiction. \\
    
    Let us assume that there is a sequence of $\Upsilon'_n \in \mathrm{Teich}(S_g^b)$ such that 
    \begin{itemize}
        \item there exists $\Upsilon_n \in \mathcal{U}_\chi^S$ with $\mathcal{O}_S(\Upsilon'_n)=\mathcal{O}_S(\Upsilon_n)$ and\\
        \item such that any number of boundary components different from $\gamma_{j_0}$ has its length that goes to $0$ as $n$ goes to infinity.
    \end{itemize}
    Set $0<k<b-2$ and let $\gamma_{j_{k+1}},...,\gamma_{j_{b-1}}$ be the boundary components whose length goes to $0$ as $n$ goes to infinity. \\
    
    %%%%%%%%%%%%%
    For each $n$, let $\varepsilon_n>0$ as in Proposition~\ref{prop:Hummel}. Thus, for all $\varepsilon'>0$, up to a subsequence, the $thick_{\varepsilon_n}(\Upsilon'_n)$ are all \\
    $(1+\varepsilon')$-bilipschitz. Fix $\varepsilon'$, then set
    \[
        N(\Upsilon'_n,L)= \# \{ \tau \mbox{ simple orthogeodesic in } \Upsilon'_n \mid \ell_{\Upsilon'_n}(\tau) < L \}
    \]
    and
    \[
        M(\Upsilon'_n,L) = \# \{ \tau \mbox{ simple orthogeodesic in } thick_{\varepsilon_n}(\Upsilon'_n) \mid \ell_{\Upsilon'_n}(\tau) < L \}.
    \]

    
    By Theorem~\ref{thm:comptageArc}, there exists $L_{\varepsilon'}$ big enough such that
    \begin{align}
        N \left( \Upsilon'_0,\frac{L_{\varepsilon'}}{q_\chi^S(1+\varepsilon')} \right) > & (1-\varepsilon')A'\mathfrak{m}(\Upsilon'_0)\left( \frac{L_{\varepsilon'}}{q_\chi^S(1+\varepsilon')} \right)^{6g-6+2b} & \mbox{ and}\label{rel:comptage1generic} \\
        M \left( \Upsilon'_0,L_{\varepsilon'} \right) < & (1+\varepsilon')B'\mathfrak{m}(\Upsilon'_0) L_{\varepsilon'} ^{6g-6+2b}, \label{rel:comptage2generic}
    \end{align}
    
    For $n$ large enough, such that any arc with an endpoint on any $\gamma_{j_{k+1}},...,\gamma_{j_{b-1}}$ has length greater than $\frac{L_{\varepsilon'}}{(1+\varepsilon')} $, we have 
    \[
        N \left( \Upsilon'_n, \frac{L_{\varepsilon'}}{(1+\varepsilon')} \right) = M\left( \Upsilon'_n,\frac{L_{\varepsilon'}}{(1+\varepsilon')} \right).
    \]
    As $\mathcal{O}_S(\Upsilon'_n)=\mathcal{O}_S(\Upsilon_n)$ for all $n$ we have $N \left( \Upsilon'_n, \frac{L_{\varepsilon'}}{(1+\varepsilon')} \right) = N \left( \Upsilon_n, \frac{L_{\varepsilon'}}{(1+\varepsilon')} \right)$. \\
    As $\Upsilon_n,\Upsilon_0 \in \mathcal{U}_\chi^S$ for all $n$, by Wolpert's~Lemma~\ref{lemme:OrthoWolpert} we have 
    \[
        N \left( \Upsilon_n, \frac{L_{\varepsilon'}}{(1+\varepsilon')} \right) \geqslant N \left( \Upsilon_0, \frac{L_{\varepsilon'}}{q_\chi^S(1+\varepsilon')} \right) = N \left( \Upsilon'_0, \frac{L_{\varepsilon'}}{q_\chi^S(1+\varepsilon')} \right).
    \]
    Moreover, $\Upsilon'_n$ and $\Upsilon'_0$ are $(1+\varepsilon')$-bilipschitz thus $M\left( \Upsilon'_n,\frac{L_{\varepsilon'}}{(1+\varepsilon')} \right) \leqslant M(\Upsilon'_0,L_{\varepsilon'})$. We obtain
    \[
        N \left( \Upsilon'_0,\frac{L_{\varepsilon'}}{q_\chi^S(1+\varepsilon')} \right) \leqslant M(\Upsilon'_0,L_{\varepsilon'}).
    \]
    By (\ref{rel:comptage1generic}) and (\ref{rel:comptage2generic}) we have
    \[
        (1-\varepsilon')A'\mathfrak{m}(\Upsilon'_0)\left( \frac{L_{\varepsilon'}}{q_\chi^S(1+\varepsilon')} \right)^{6g-6+2b} < (1+\varepsilon')B'\mathfrak{m}(\Upsilon'_0) L_{\varepsilon'}^{6g-6+2b}.
    \]
    Thus, for all $\varepsilon'>0$ we have,
    \[
        A' \left ( \frac{1}{q_\chi^S} \right)^{6g+2b-6}(1-\varepsilon') < B'(1+\varepsilon')^{5-6g-2b}.
    \]
    At the limit, when $\varepsilon'$ goes to $0$ we obtain that $A'\left ( \frac{1}{q_\chi^S} \right)^{6g+2b-6} \leqslant B'$ which is a contradiction. Thus, there exists a lower bound $\varepsilon_\chi^3$ on the length of each boundary component. 
    
    
    %%%%%%%%%%%%%%%%%%%%%
    Set $\varepsilon_\chi^S=\min(\varepsilon_\chi^1,\varepsilon_\chi^2, \varepsilon_\chi^3)$. Then by Corollary~\ref{cor:mumf}, there exists a surface isometric to $\Upsilon'$ in $\mathcal{Q}(A_\chi^S,\varepsilon_\chi^S)$.
\end{proof}
%%%%%%%%%%%%%%%%%%%%%%%%%%%%%%%%%%%%%%%%%%%%%


%%%%%%%%%%%%%%


Fix a hexagon decomposition $\mathcal{H}=\{\alpha_1,...,\alpha_{6g+3b-6} \}$ on $S_g^b$. Let $\chi_0 \in \mathrm{Teich}(S_g^b)$ be given by 
\[
    (\ell_{\chi_0}(\alpha_1),...,\ell_{\chi_0}(\alpha_{6g+3b-6}))=(1,...,1).
\]
This surface is going to be a point of reference. Recall that, for any \mbox{$\omega=(\ell_1,...,\ell_{6g+3b-6})\in \mathbb{R}_+^{6g+3b-6}$}, we set $\chi^{\omega}:=\varphi_{\mathcal{H}}^{-1}(\omega) \in \mathrm{Teich}(S_g^b)$, see Definition~\ref{def:omegaSurface}. Then we choose a quasi-conformal homeomorphism~$\phi^{\omega}$ between $\chi_0$ and $\chi^{\omega}$. For each orthogeodesic $\tau$ on $\chi_0$, we denote by $\tau(\chi^{\omega})$ the unique orthogeodesic in the free homotopy class of the orthogeodesic $\phi^{\omega} \circ (\tau) $ on $\chi^{\omega} $. For any finite or infinite ordered set $\mathbb{O}$ of orthogeodesics $\tau_1, \tau_2, ...$ on $\chi_0$, we define the sequences
\begin{align*}
    \mathbb{O}(\Upsilon) & =(\tau_1(\Upsilon), \tau_2(\Upsilon),...) \\
    \ell_{\Upsilon}(\mathbb{O}) & =(\ell_{\Upsilon}(\tau_1), \ell_{\Upsilon}(\tau_2), ... ).
\end{align*}

We set $\Pi= (\beta_1, \beta_2, ...)$ to be the sequence of all orthogeodesics on $\chi_0$, arranged so that $\ell_{\chi_0}(\beta_1) \leqslant \ell_{\chi_0}(\beta_2) \leqslant ...$ and set $\Pi_k = (\beta_1,...,\beta_k)$. Then, set $\Pi_k^S= (\beta_1^S, \beta_2^S, ...)$ the sequence of all simple orthogeodesics on $\chi_0$, arranged so that $\ell_{\chi_0}(\beta_1^S) \leqslant \ell_{\chi_0}(\beta_2^S) \leqslant ...$ and set $\Pi_k^S=(\beta_1^S,...,\beta_k^S)$. Note that $\ell_{\chi_0}(\Pi)=\mathcal{O}(\chi_0)$ and $\ell_{\chi_0}(\Pi^S)=\mathcal{O}_{S}(\chi_0)$. Finally, let $\delta_1,..,\delta_{6g+3b-6}$ be the unique collection of simple orthogeodesics on $\chi_0$ such that $i(\alpha_i,\delta_j)=\delta_{ij}$ for all $1 \leqslant i,j \leqslant 6g+3b-6$, and set $\Sigma = \mathcal{H} \cup \{ \delta_1,...,\delta_{6g+3b-6} \}$.\\


%%%%%%%%%%%%%%%%%

We fix $\chi \in \mathrm{Teich}(S_g^b)$. In the case of the orthospectrum, we assume that $\chi \in \mathrm{Teich}(S_g^b) \setminus \bigcup_{[[\tau]]\neq [[\tau']]} \mathcal{E}([[\tau]],[[\tau']])$ but we do not impose such a restriction in the case of the simple orthospectrum. Let $\mathcal{U}_\chi$ and $A_\chi,\varepsilon_\chi>0$ be as in Proposition~\ref{prop:Ux}. Fix $q_0^S >1$, let $\mathcal{U}_\chi^S$ be an open ball of center $\chi$ and radius $\frac{\log(q_0^S)}{2}$, and $A_\chi^S,\varepsilon_\chi^S>0$ be as in Proposition~\ref{prop:UxS}. The open neighborhood $\mathcal{U}_\chi$ has a compact closure and there is a compact $\mathcal{Q}(A_\chi,\varepsilon_\chi)$ that contains it. Similarly, the open neighborhood $\mathcal{U}_\chi^S$ has a compact closure and there is a compact $\mathcal{Q}(A_\chi^S,\varepsilon_\chi^S)$ that contains it.\\

%%%%%%%%%%%%%%%

Let $D = \{ \chi^{\omega} \in \mathrm{Teich}(S_g^b) \mid  \frac{1}{10} \leqslant \ell _{\chi^{\omega}}(\alpha_1),...,\ell_{\chi^{\omega}} (\alpha_{6g+3b-6}) \leqslant 10 \} $. In the case of the orthospectrum, let $C, C_1$ be compact sets of $\mathrm{Teich}(S_g^b)$ whose interiors $\mathring{C}, \mathring{C_1} \subset \mathrm{Teich}(S_g^b)$ are connected and
\[
    ( \mathcal{Q}(A_\chi, \varepsilon_\chi) \cup D ) \subset \mathring{C} \subset C \subset \mathring{C_1}.
\]
And for the simple orthospectrum, let $C^S, C_1^S$ be compact sets whose interiors $\mathring{C^S}, \mathring{C_1^S}$ are connected and
\[
    ( \mathcal{Q}(A_\chi^S, \varepsilon_\chi^S) \cup D ) \subset \mathring{C^S} \subset C^S \subset \mathring{C_1^S}.
\]
By Corollary~\ref{cor:Wolp}, there exists $q \geqslant 1$ such that 
\begin{align}
   \frac{\ell_{\Upsilon} (\beta)}{q}  \leqslant \ell_{\Upsilon'} (\beta) \leqslant q \ell_{\Upsilon} (\beta)\label{rel:q} 
\end{align}
for any $\Upsilon, \Upsilon' \in C_1$ and any $\beta \in \Pi$. This $q$ will remain fixed during the proof. We denote $q^S$ the constant obtained for any $\Upsilon, \Upsilon' \in C_1^S$.


\begin{lemme}\label{lem:rho}
    For any $k \in \mathbb{N} $, there exists an integer $k^* \geqslant k $, which depends only on $k$ and $C_1$, with the following property. If $\Upsilon', \Upsilon'' \in C_1$ and if $\rho : \Pi_k \to \Pi$ is an injection such that
    \[
        \ell_{\Upsilon'} (\Pi_k) = \ell_{\Upsilon''}(\rho (\Pi_k))
    \]
    then $\rho(\Pi_k) \subset \Pi_{k^*}$.
    The same is true with $\Pi^S, \Pi_k^S$, $\Pi_{k^*}^S$ and $C_1^S$ instead of $\Pi, \Pi_k$, $\Pi_{k^*}$ and $C_1$.
\end{lemme}


\begin{proof}
    Let $c_1=\max\{ \ell_{\Upsilon} (\beta) \mid  \Upsilon \in C_1, \beta \in \Pi_k \} $, and let $k^* \geqslant k$ be such that, on the base surface $\chi_0$, we have $\ell_{\chi_0} (\beta_j) > q c_1$ for all $j > k^*$. By (\ref{rel:q}) and since $\chi_0$ is in $D \subset C_1$, we have $ \ell_{\Upsilon} (\beta_j) \geqslant \frac{\ell_{\chi_0} (\beta_j)}{q} > c_1 $ for all $j>k^*$ and $\Upsilon \in C_1$. Since we have $\ell_{\Upsilon''} (\rho (\beta_i))=\ell_{\Upsilon'} (\beta_i) \leqslant c_1$ for $ i \leqslant k$ by definition of $c_1$, it follows that $\rho (\beta_i) \in \Pi_{k^*}$. For the simple orthospectrum, just replace $\beta_i$, $C_1$, $\Pi, \Pi_k$ and $\Pi_{k^*}$ with $\beta_i^S$, $C_1^S$,  $\Pi^S, \Pi_k^S$ and $\Pi_{k^*}^S$. 
\end{proof}











\subsection{The first lengths of the orthospectrum}

In this section, we will see that any pair $\Upsilon$, $\Upsilon'$ in a compact subset of the Teichmüller space have the same orthospectrum, or the same simple orthospectrum if and only if they share a certain finite number of length in their orthospectrum or simple orthospectrum. This is a key step of the proof of Theorem~\ref{thm:Main2}.\\

We define for each $k \in \mathbb{N}$ the following sets:
\begin{align*}
    V_k^1 & =\{ (\Upsilon, \Upsilon') \in \mathring{C_1} \times \mathring{C_1} \mid  \ell_{\Upsilon} (\Pi_k) \subset \mathcal{O}(\Upsilon') \text{ and } \ell_{\Upsilon'} (\Pi_k) \subset \mathcal{O}(\Upsilon) \} \\
    V_k & = V_k^1 \cap (\mathring{C} \times \mathring{C}) \\
    W_k^1 & =\{ (\Upsilon, \Upsilon') \in \mathring{C_1^S} \times \mathring{C_1^S} \mid  \ell_{\Upsilon} (\Pi_k^S) \subset \mathcal{O}_S(\Upsilon') \text{ and } \ell_{\Upsilon'} (\Pi_k^S) \subset \mathcal{O}_S(\Upsilon) \} \\
    W_k & = W_k^1 \cap (\mathring{C^S} \times \mathring{C^S}).
\end{align*}
Let $k^*, k_S^*\in \mathbb{N} $ be as in Lemma~\ref{lem:rho}. Given two pairs of injections \mbox{$\rho_1, \rho_2 : \Pi_k \to \Pi_{k^*}$} and $\rho_1^S, \rho_2^S : \Pi_k^S \to \Pi_{k_S^*}^S$, we set
\begin{align*}
    V[\rho_1,\rho_2]=\{ (\Upsilon, \Upsilon') \in \mathring{C_1} \times \mathring{C_1} \mid  \ell_{\Upsilon} (\Pi_k) = \ell_{\Upsilon'} (\rho_1 (\Pi_k)) \text{ and } \ell_{\Upsilon'}( \Pi_k ) = \ell_{\Upsilon} (\rho_2 (\Pi_k)) \} \\
    W[\rho_1^S,\rho_2^S]=\{ (\Upsilon, \Upsilon') \in \mathring{C_1^S} \times \mathring{C_1^S} \mid  \ell_{\Upsilon} ( \Pi_k^S ) = \ell_{\Upsilon'} (\rho_1^S (\Pi_k^S)) \text{ and } \ell_{\Upsilon'} (\Pi_k^S ) = \ell_{\Upsilon} (\rho_2^S (\Pi_k^S)) \}.
\end{align*}
By Lemma~\ref{lem:RealAnalLength}, the spaces $V[\rho_1,\rho_2]$ and $W[\rho_1^S,\rho_2^S]$ are real analytic subvarieties of \\
$\mathring{C_1} \times \mathring{C_1}$ and $\mathring{C_1^S} \times \mathring{C_1^S}$. Then, thanks to Lemma~\ref{lem:rho}, we have
\begin{align*}
    V_k^1= \bigcup_{(\rho_1, \rho_2)} V[\rho_1,\rho_2] \\
    W_k^1= \bigcup_{(\rho_1^S, \rho_2^S)} W[\rho_1^S,\rho_2^S].
\end{align*}
Since there are finitely many pairs $(\rho_1, \rho_2)$ and $(\rho_1^S,\rho_2^S)$, the unions $V_k^1$ and $W_k^1$ are also real analytic varieties.\\

    
    
\begin{lemme}\label{lem:K}
    There exists $K \in \mathbb{N}$ such that $V_{K+j}=V_K$ and $W_{K+j}=W_K$ for all $j \geqslant 1$.
\end{lemme}

\begin{proof}
    The Teichmüller space is a real analytic space, and $C_1$ and $C_1^S$ are a compact subsets of $\mathrm{Teich}(S_g^b)$. Thus, by \cite[Théorème I.9]{JFrisch}, the rings of real analytic functions on~$C_1 \times C_1$ and~$C_1^S \times C_1^S$ are Noetherian. Moreover, any real analytic subvariety $V$ is associated with the ideal $I(V)$ of real analytic functions vanishing on the subvariety. For two real analytic subvarieties $V_1, V_2$, the inclusion $V_1 \subset V_2$ is equivalent to \\
    $I(V_2) \subset I(V_1)$ (see \cite{SousVar} for more details). By definition, increasing sequences of ideals of a Noetherian ring are stationary, so decreasing sequences of subvarieties are stationary. Therefore, there is $K \in \mathbb{N}$ such that $V_{K+j}=V_K$ and $W_{K+j}=W_K$ for all $j \geqslant 1$. 
\end{proof}
    
    
    








\subsection{Proof of Theorem~\ref{thm:Main2}}

We define $\mathcal{U}_\chi \subset \mathrm{Teich}(S_g^b), D, C$ and $C_1$ as in Section~\ref{subsec:SetUp}. To avoid repetition, we prove Theorem~\ref{thm:Main2} only for the orthospectrum. For the simple orthospectrum, the proof is the mostly the same with $\Pi^S$, $\mathcal{U}_\chi^S$, $C^S$ and $C_1^S$ instead of $\Pi$, $\mathcal{U}_\chi$, $C$ and $C_1$ and so on. We will note when there are variations between the two proofs, for example some verification we perform about the simplicity of curves are not necessary when proving the theorem for the simple orthospectrum.

\begin{proof}[Proof of Theorem~\ref{thm:Main2}]
    We fix $K$ as in Lemma~\ref{lem:K} and large enough so that $\Sigma \subset \Pi_K$. Then, with the notations of Lemma~\ref{lem:rho}, we fix $M=K^*$ and $N=M^*$. 
    Now, \\
    let $\rho : \Pi_M \to \Pi_N$ be any injection such that $\Pi_K \subset \rho (\Pi_M)$. We define
    \begin{align*}
        V_{\rho} & =\{\Upsilon \in \mathring{C} \mid\text{ there exists } \Upsilon^{\rho} \in \mathrm{Teich}(S_g^b) \text{ such that } \ell_{\Upsilon} (\Pi_M) = \ell_{\Upsilon^{\rho}} ( \rho (\Pi_M))\},\\
        \overline{V_{\rho}} &  =\{\Upsilon \in \mathring{C} \mid\text{ there exists } \Upsilon^{\rho} \in \mathring{C} \text{ such that } \ell_{\Upsilon} (\Pi_M) = \ell_{\Upsilon^{\rho}} ( \rho (\Pi_M))\}.
    \end{align*}
    For $\Upsilon \in V_{\rho} $, the surface $\Upsilon^{\rho} $ is unique: indeed, since \mbox{$\mathcal{H} \subset \Sigma \subset \Pi_K \subset \rho (\Pi_M)$}, the vector $\ell_{\Upsilon^{\rho}} (\mathcal{H} )= \ell_{\Upsilon} (\rho^{-1}(\mathcal{H}))  $ represents the Ushijima coordinates of $\Upsilon^{\rho}$. Moreover, we have $\overline{V_{\rho}} \subset V_{\rho}$, and for $\Upsilon \in \overline{V_{\rho}} $, we have $\mathcal{O}(\Upsilon)=\mathcal{O}(\Upsilon^{\rho} ) $. We have
    \[
        \ell_{\Upsilon^{\rho}} (\Pi_K ) \subset \ell_{\Upsilon^{\rho}} (\rho (\Pi_M))=\ell_{\Upsilon} (\Pi_M) \subset \mathcal{O}(\Upsilon) .
    \]
    Conversely, since $\Pi_K \subset \Pi_M$, we also have $\ell_{\Upsilon} (\Pi_K) \subset \ell_{\Upsilon} (\Pi_M) = \ell_{\Upsilon^{\rho}} (\rho(\Pi_M)) \subset \mathcal{O}(\Upsilon^{\rho}) $ so $\mathcal{O}(\Upsilon)=\mathcal{O}(\Upsilon^{\rho})$ by Lemma~\ref{lem:K}. \\
    
    Now, we define the real analytic mapping $m_{\rho} : \mathrm{Teich}(S_g^b) \to \mathrm{Teich}(S_g^b)$ given by $m_{\rho}(\chi^{\omega})= \chi^{\ell_{\chi^{\omega}} (\rho^{-1}(\mathcal{H}))}$. If $\Upsilon^{\rho}$ exists, we have $m_{\rho}(\Upsilon)=\Upsilon^{\rho}$. Thus, 
    \[
        V_{\rho}=\{\Upsilon \in \mathring{C} \mid \ell_{\Upsilon} (\Pi_M) = \ell_{m_{\rho}(\Upsilon)} ( \rho (\Pi_M)) \}
    \]
    is a real analytic subvariety of $\mathring{C}$.\\
    We define 
    \begin{align*}
        V_{\rho}^*=\{ \Upsilon \in V_{\rho} \mid \Upsilon^{\rho} \text{ is isometric to } \Upsilon \},\\
        \overline{V_{\rho}^*}=\{ \Upsilon \in \overline{V_{\rho}} \mid \Upsilon^{\rho} \text{ is isometric to } \Upsilon \}.
    \end{align*}
    A surface $\Upsilon \in V_{\rho}$ is isometric to $\Upsilon^{\rho}=m_{\rho}(\Upsilon)$ if and only if $\Upsilon$ and $\Upsilon^{\rho}$ are in the same $\mathcal{MCG}(S_g^b)$ orbit, that is, if and only if there is an injection $r : \mathcal{H} \to \Pi$ induced by a homeomorphism of the base surface $\chi_0$ such that 
    \[
        \ell_{\Upsilon} (\mathcal{H})=\ell_{m_{\rho}(\Upsilon)} (r (\mathcal{H})).
    \]
    By Lemma~\ref{lem:rho}, we have $r (\mathcal{H}) \subset \Pi_M$. This shows that the set $R^*$ of all such possible injections $r$ is finite. This implies that
    \[
        V_{\rho}^*= \bigcup_{r \in R^*}\{\Upsilon \in V_{\rho} \mid \ell_{\Upsilon} (\mathcal{H})=\ell_{m_{\rho}(\Upsilon)} (r (\mathcal{H}))\}
    \]
    is a real analytic subvariety of $V_{\rho}$.

    Let $\mathcal{R}$ be the set of injective maps $\rho : \Pi_M \to \Pi_N$ satisfying $\Pi_K \subset \rho (\Pi_M)$. Note that this set is finite. Define
    \begin{align*}
        V = \bigcup_{\rho \in \mathcal{R}}(V_{\rho} \setminus V_{\rho}^*),\\
        \overline{V} = \bigcup_{\rho \in \mathcal{R}}(\overline{V_{\rho}} \setminus \overline{V_{\rho}^*}).
    \end{align*}
    The sets $V_{\rho}, V_{\rho}^*$ are real analytic subvarieties so $V$ is a real analytic subvariety of $\mathring{C}$. Note that $\overline{V_{\rho}^*} \subset V_{\rho}^*$, thus $(\overline{V_{\rho}} \setminus \overline{V_{\rho}^*}) \subset (V_{\rho} \setminus V_{\rho}^*)$ and $\overline{V}$ is included in a real analytic subvariety of $\mathring{C}$. \\
    
    
    
    
    By construction, we have $\overline{V} \cap \mathring{C} \subset \mathcal{V}_g^b \cap \mathring{C}$, hence $\overline{V} \cap \mathcal{U}_\chi \subset \mathcal{V}_g^b \cap \mathcal{U}_\chi$. Conversely, if $\Upsilon \in \mathcal{V}_g^b \cap \mathcal{U}_\chi$ then there exists $\Upsilon'$ not isometric to $\Upsilon$ such that $\mathcal{O}(\Upsilon)=\mathcal{O}(\Upsilon')$. By Proposition~\ref{prop:Ux}, we can assume that $\Upsilon' \in \mathcal{Q}(A_\chi,\varepsilon_\chi) \subset \mathring{C}$. Moreover, there exists a bijection $\rho : \Pi \to \Pi$ satisfying $\ell_{\Upsilon} (\Pi)=\ell_{\Upsilon'}(\rho (\Pi)) $. By Lemma~\ref{lem:rho}, we have $\rho^{-1}(\Pi_K) \subset \Pi_M $ and $\rho(\Pi_M) \subset \Pi_N$. In other words, there exists an injection $\rho : \Pi_M \to \Pi_N$ such that
    \begin{align*}
        &\Pi_K \subset \rho(\Pi_M) \\
        &\ell_{\Upsilon} (\Pi_M) = \ell_{\Upsilon'} (\rho (\Pi_M)).
    \end{align*}
    By definition of $\overline{V}$, we have $\Upsilon \in V \cap \mathcal{U}_\chi$. In conclusion, $\overline{V} \cap \mathcal{U}_\chi = \mathcal{V}_g^b \cap \mathcal{U}_\chi \subset V \cap \mathcal{U}_\chi$ for any neighborhood $\mathcal{U}_\chi$. \\
    
    Let us show that $\overline{V} \cap \mathcal{U}_\chi = \mathcal{V}_g^b \cap \mathcal{U}_\chi \neq \mathcal{U}_\chi$. We proceed by showing that $V \cap \mathcal{U}_\chi \neq \mathcal{U}_\chi$, \textit{i.e.}, that $\dim V < \dim \mathrm{Teich}(S_g^b)$. Since $V_{\rho}$ is a real analytic subvariety of $\mathring{C}$ and $\mathring{C}$ is connected by definition, we either have $V_{\rho}=\mathring{C}$ or else $\dim V_{\rho} < \dim \mathring{C} = \dim \mathrm{Teich}(S_g^b)$. We want to show that if $V_{\rho}=\mathring{C}$ then $V_{\rho}=V_{\rho}^*$. If this is true, for all $\rho \in \mathcal{R}$, \\
    $\dim(V_{\rho}\setminus V_{\rho}^*)<\dim \mathrm{Teich}(S_g^b)  $, thus $\dim V <\dim \mathrm{Teich}(S_g^b)$.\\
    
    So suppose $V_{\rho}=\mathring{C}$, then there is a map $m_{\rho} : \mathring{C} \to \mathrm{Teich}(S_g^b)$ such that \\
    $\ell_{m_{\rho}(\Upsilon)} (\rho (\Sigma))=\ell_{\Upsilon} (\Sigma)$ for all $\Upsilon \in \mathring{C}$. In the following, we abbreviate $\Tilde{\beta}:=\rho(\beta)$ for any $\beta \in \Sigma$, and $\Tilde{\Upsilon}:=m(\Upsilon)$ for any $\Upsilon \in \mathring{C}$.
    
    \underline{Step 1:} $\rho(\mathcal{H})$ is a hexagon decomposition.\\
    Set $\omega=(\frac{1}{4},...,\frac{1}{4})$. Then $\ell_{\Tilde{\chi}^{\omega}} (\Tilde{\alpha}_i)=\frac{1}{4}$ for $1 \leqslant i \leqslant 6g+3b-6$. By Theorem~\ref{thm:NonSimpleLength}, we deduce that the orthogeodesics $\Tilde{\alpha}_i$ are simple. Note that in the case of the simple orthospectrum we already know that these orthogeodesics are simple. The fact that they are pairwise disjoint follows from Lemma~\ref{lemme:orthoCollier}. If $\Tilde{\alpha}_i$ and $\Tilde{\alpha}_j$ intersect each other, then \mbox{$\sinh(\ell_{\Tilde{\chi}^{\omega}} (\Tilde{\alpha}_i))\sinh (\ell_{\Tilde{\chi}^{\omega}} (\Tilde{\alpha}_j)) > 1$}. This is impossible since $\sinh(\frac{1}{4})\sinh(\frac{1}{4}) <1$. Therefore, $\rho$ sends $\mathcal{H}$ to another hexagon decomposition $\Tilde{\mathcal{H}}$ of $\chi_0$. 

    \underline{Step 2:} Understand the relative position of the $\Tilde{\alpha}_i$.\\
    We want to show that if $\alpha_{i_1}, \alpha_{i_2}, \alpha_{i_3}, \alpha_{i_4}$ are orthogeodesics delimiting an octagon of orthogonals $\alpha_j$ and $\delta_j$ on $\chi_0$, then $\Tilde{\alpha}_{i_1}, \Tilde{\alpha}_{i_2}, \Tilde{\alpha}_{i_3}, \Tilde{\alpha}_{i_4}$ are also orthogeodesics delimiting an octagon of orthogonals $\Tilde{\alpha}_j$ and $\Tilde{\delta}_j$. Indeed, fix $\Upsilon$ and $1 \leqslant j \leqslant 6g+3b-6$ such that $\ell_{\Upsilon} (\alpha_i) = \frac{1}{4}$ and $\ell_{\Upsilon} (\alpha_j) = 6$ for all $1 \leqslant i\neq j \leqslant 6g+3b-6$. By Lemma~\ref{lem:OctaDroit}, we have
    \[
        \cosh(\ell_{\Upsilon} (\delta_j)) =f_{\mathrm{octa}}(\ell_{\Upsilon} (\alpha_j), \ell_{\Upsilon} (\alpha_{i_1}),\ell_{\Upsilon} (\alpha_{i_2}),\ell_{\Upsilon} (\alpha_{i_3}), \ell_{\Upsilon} (\alpha_{i_4} ))
    \]
    and $\ell_{\Tilde{\Upsilon}} (\Tilde{\delta}_j)=\ell_{\Upsilon} (\delta_j) < \frac{1}{2}$. As before, by Theorem~\ref{thm:NonSimpleLength} and Lemma~\ref{lemme:orthoCollier}, the curve~$\Tilde{\delta}_j$ is simple and disjoint from $\Tilde{\alpha}_i$ for any $1 \leqslant i\neq j \leqslant 6g+3b-6$. It follows that $i(\Tilde{\delta}_j,\Tilde{\alpha}_j)=1$ and the arcs $\Tilde{\delta}_j, \Tilde{\alpha}_j$ lie in an octagon delimited by four orthogeodesics among the $\Tilde{\alpha}_i$ for $i \neq j$. Since
    \[
        \cosh(\ell_{\Tilde{\Upsilon}} (\Tilde{\delta}_j)) =f_{\mathrm{octa}}(\ell_{\Tilde{\Upsilon}} (\Tilde{\alpha}_j), \ell_{\Tilde{\Upsilon}} (\Tilde{\alpha}_{i_1}),\ell_{\Tilde{\Upsilon}} (\Tilde{\alpha}_{i_2}),\ell_{\Tilde{\Upsilon}} (\Tilde{\alpha}_{i_3}), \ell_{\Tilde{\Upsilon}} (\Tilde{\alpha}_{i_4}) ),
    \]
    we see that if we vary the length of one $\Tilde{\alpha}_l \in \Tilde{\mathcal{H}}$ and fix the other ones, the length of $\Tilde{\delta}_j$ only depends on the lengths of $\Tilde{\alpha}_{i_1},\Tilde{\alpha}_{i_2},\Tilde{\alpha}_{i_3},\Tilde{\alpha}_{i_4}$ and $\Tilde{\alpha}_j$. This means that the orthogeodesics delimiting the octagon containing $\Tilde{\delta}_j$ and $\Tilde{\alpha_j}$ are $\Tilde{\alpha}_{i_1},\Tilde{\alpha}_{i_2},\Tilde{\alpha}_{i_3},\Tilde{\alpha}_{i_4}$. 
    
    Moreover, the non-symmetry of $f_{\mathrm{octa}}$ also gives an indication about how the orthogeodesics delimiting the octagon are placed. Indeed, if $\alpha, \delta_1, \delta_2, \delta_3$ and $\delta_4$ are as in Figure~\ref{fig:DemoOctagon} then we obtain different values of $\beta$ when we exchange $\delta_1$ and $\delta_2$ or when we exchange $\delta_1$ and $\delta_4$.  
    Thus, if we choose $\Upsilon$ such that $\alpha_{i_1}$,$\alpha_{i_2},\alpha_{i_3},\alpha_{i_4},\alpha_{j}$ and~$\delta_j$ all have different lengths. Then, we know that if $(\alpha_{i_1}, \alpha_{i_2},\alpha_j)$ delimits a hexagon, then $(\Tilde{\alpha}_{i_1}, \Tilde{\alpha}_{i_2},\Tilde{\alpha}_j)$ also delimits an hexagon. Thus, the two surfaces are isometric because they are obtained by gluing isometric hexagons with the same pattern. \\
    In conclusion, if $V_{\rho}=\mathring{C}$, then $V_{\rho}=V_{\rho}^*$ and we have $V_{\rho} \setminus V_{\rho}^* = \emptyset$ and $\overline{V} \cap \mathcal{U}_\chi = \mathcal{V}_g^b \cap \mathcal{U}_\chi \neq \mathcal{U}_\chi$. \\
    
    In the case of the simple orthospectrum, the union of the $\mathcal{U}_\chi^S$ covers the whole Teichmüller space, thus it is enough to say that $\mathcal{W}_g^b$ is a proper local real analytic subvariety of $\mathrm{Teich}(S_g^b)$ thus a negligible set. In the case of the orthospectrum, the union of the $\mathcal{U}_\chi$ does not necessarily cover the whole Teichmüller space. However, it covers $\mathrm{Teich}(S_g^b) \setminus \bigcup_{[[\tau]]\neq [[\tau']]} \mathcal{E}([[\tau]],[[\tau']])$, which is dense in the Teichmüller space, thus $\mathcal{V}_g^b \cap ( \mathrm{Teich}(S_g^b) \setminus \bigcup_{[[\tau]]\neq [[\tau']]} \mathcal{E}([[\tau]],[[\tau']]) )$ is a local proper real analytic variety and is a negligible set of $\mathrm{Teich}(S_g^b)$. By Proposition~\ref{prop:E_negligible}, $\bigcup_{[[\tau]]\neq [[\tau']]} \mathcal{E}([[\tau]],[[\tau']])$ is also a negligible set of $\mathrm{Teich}(S_g^b)$. Thus their union is a negligible set and
    \[
        \mathcal{V}_g^b \subset \left(  \mathcal{V}_g^b \cap ( \mathrm{Teich}(S_g^b) \setminus \bigcup_{[[\tau]]\neq [[\tau']]} \mathcal{E}([[\tau]],[[\tau']]) )  \right) \bigcup_{[[\tau]]\neq [[\tau']]} \mathcal{E}([[\tau]],[[\tau']]).
    \]
    is included in a negligible set of $\mathrm{Teich}(S_g^b)$.
\end{proof}












































\section{Rigidity results}\label{sec:4}

In \cite{orthoSys}, Masai and McShane prove orthospectrum rigidity for the one-holed torus. In the case of the simple orthospectrum, the same proof does not work as it relies on computing the length of the boundary using Basmajian's identity. However in this section, we prove simple orthospectrum rigidity for the one-holed torus with a different proof, which relies on Ushijima coordinates instead of Fenchel-Nielsen coordinates. 

The first result we need to prove rigidity is the following:

\begin{proposition}\label{prop:DisjFirst}
    Let $X$ be a compact hyperbolic surface with geodesic boundary. Then the first two lengths in $\mathcal{O}_S(X)$ are the lengths of two disjoint orthogeodesics.
\end{proposition}

\begin{proof}
    Let $\tau_1$ and $\tau_2$ be two orthogeodesics realizing the first two lengths of the simple orthospectrum with $\ell(\tau_1) \leqslant \ell(\tau_2)$. Let us suppose that $i(\tau_1,\tau_2)>0$. The idea is to get a contradiction by constructing a new orthogeodesic $\tau$ shorter than $\tau_2$.

    We construct a piecewise geodesic path $\sigma$ as follows. Start at an endpoint of $\tau_1$ such that the length between this endpoint and the first intersection point $p$ between $\tau_1$ and $\tau_2$ is less than $\frac{\ell(\tau_1)}{2}$. Then follow $\tau_1$ until $p$, and finally follow $\tau_2$ until its closest endpoint. We obtain $\ell(\sigma) \leqslant \frac{\ell(\tau_2)}{2} + \frac{\ell(\tau_1)}{2} \leqslant \ell(\tau_2)$. Note that $\sigma$ is essential, otherwise, together with an arc of the boundary of $X$, we get a hyperbolic triangle with two right angles, which is impossible.

    \begin{figure}[H]
        \centering
        \includegraphics[height=7cm]{figures/intersecte.png}
        \caption{Construction of $\sigma$.}
        \label{fig:Inter}
    \end{figure}

    Note that the orthogeodesic $\tau$ homotopic to $\sigma$ is simple. Since the two arcs forming~$\sigma$ meet at some angle different from~$\pi$, the geodesic representative~$\tau$ is strictly shorter than $\sigma$ and $\ell(\tau) < \ell(\sigma) \leqslant \ell(\tau_2)$. Moreover, by construction \mbox{$i(\tau,\tau_1)<i(\tau_1,\tau_2)$}. Suppose then that~$\tau=\tau_1$. In this case, the segment of~$\tau_2$ that~$\sigma$ follows and the segment of~$\tau_1$ that~$\sigma$ does not follow are homotopic and form with a segment of~$\partial X$ a hyperbolic triangle with two right angles. This is impossible, so $\tau \neq \tau_1$. 

    Thus, the simple orthogeodesic $\tau$ is different from $\tau_1$ and shorter than $\tau_2$, which is a contradiction.
\end{proof}

With this result at hand, we can prove the desired rigidity statement.

\begin{theorem}\label{thm:Main3}
    Let $T$ and $T'$ be two hyperbolic structures with geodesic boundary on the one-holed torus.
    Then $T$ and $T'$ are isometric if and only if $\mathcal{O}_S(T)=\mathcal{O}_S(T')$.
\end{theorem}



\begin{proof}
    A hexagon decomposition of a one-holed torus is formed by three arcs. Our goal is to find a hexagon decomposition where the three orthogeodesics have length \\
    $t_1 \leqslant t_2 \leqslant t_3$, which are the first three lengths of the simple orthospectrum.

    By Proposition~\ref{prop:DisjFirst}, the first two lengths $t_1$ and $t_2$ of $\mathcal{O}_S(T)=\mathcal{O}_S(T')$ correspond to two disjoint simple orthogeodesics $\tau_1$ and $\tau_2$ on $T$ (respectively $\tau'_1$ and $\tau'_2$ on $T'$). To visualize this situation, we give $\tau_1$ and $\tau_2$ an orientation,  cut the one-holed torus along $\tau_1$ and $\tau_2$ and obtain a right-angled octagon as in Figure \ref{fig:octoTau}. 
    
    There are exactly two simple orthogeodesics, $\tau_3$ and $\Tilde{\tau}_3$, disjoint from $\tau_1$ and $\tau_2$, each of which joins two opposite sides of the octagon corresponding to arcs in $\partial X$. Assume that $\ell(\tau_3) \leqslant \ell(\Tilde{\tau_3})$.
    Our goal is to prove that any simple orthogeodesic which is not disjoint from $\tau_1$ or $\tau_2$ (or both) is longer than $\tau_3$.

    \begin{figure}[H]
        \centering
        \includegraphics[height=6cm]{figures/OctogoneTau3.png}
        \caption{$\tau_3$ and $\Tilde{\tau_3}$ on the octagon}
        \label{fig:octoTau}
    \end{figure}
    
    Let $\tau$ be a simple orthogeodesic intersecting $\tau_1$ or $\tau_2$ (or both). Without loss of generality, we can assume that $\tau$ has its endpoints on the same sides $b_1$ and $b_2$ as $\tau_3$. Indeed, if $\tau$ has its endpoints on the opposite sides, we just replace $\tau_3$ by $\Tilde{\tau_3}$ in the proof and show that $\ell(\tau_3) \leqslant \ell(\Tilde{\tau_3}) < \ell(\tau)$. Orient $\tau$ from its endpoint on $b_1$ to the endpoint on $b_2$. Assume that $\tau$ first intersects $\tau_1$ before possibly intersecting $\tau_2$ (the other case being analogous). Let $n$ be the number of time that $\tau$ intersects $\tau_1$ before possibly intersecting $\tau_2$. We prove that $\ell(\tau_3) < \ell(\tau)$ by induction on $n$.
    \\
    \\
    \underline{Base case $n=1$:} The orthogeodesic $\tau$ intersects $\tau_1$ exactly once before intersecting~$\tau_2$. 

    We denote by $p_1$ the first point of intersection between $\tau$ and $\tau_1$, and by $p_2$ the last one. Since $\tau$ is simple, $p_2$ does not lie between $p_1$ and $b_2$. In other words, we have $d(p_2,b_1)+d(p_1,b_2) \leqslant \ell(\tau_1)$. We label by $x$ the segment of $\tau$ between $p_2$ and $b_2$, and by~$y$ the segment of $\tau$ between $p_1$ and $b_1$ (as in Figure~\ref{fig:octoCas1Step1}).

    \begin{figure}[H]
        \centering
        \includegraphics[height=8.5cm]{figures/OctogoneTauCas1Etape1.png}
        \caption{Example of case $n=1$.}
        \label{fig:octoCas1Step1}
    \end{figure}

    
    We construct two arcs $\sigma$ and $\Tilde{\sigma}$ homotopic to $\tau_3$ as follows: $\sigma$ is the union of $x$ with the segment of $\tau_1$ between $p_2$ and $b_1$, and $\Tilde{\sigma}$ is the union of $y$ with the segment of $\tau_1$ between $p_1$ and $b_2$, as in Figure \ref{fig:octoCas1Step2}. We have

    \begin{align*}
        length(\sigma) = d(p_2,b_1) + length(x) \\
        length(\Tilde{\sigma}) = d(p_1,b_2) + length(y).
    \end{align*}


    \begin{figure}[H]
        \centering
        \includegraphics[height=9cm]{figures/OctogoneTauCas1Etape2.png}
        \caption{Construction of $\sigma$ and $\Tilde{\sigma}$ for case $n=1$.}
        \label{fig:octoCas1Step2}
    \end{figure}

    We obtain 
    \[
        2\ell(\tau_3) < length(\sigma) + length(\Tilde{\sigma}) \leqslant d(p_2,b_1) + d(p_1,b_2) + length(x) + length(y) \leqslant \ell(\tau_1) + \ell(\tau).
    \]
    Since $\ell(\tau_1) \leqslant \ell(\tau)$, we conclude that $\ell(\tau_3) < \ell(\tau)$.
    \\


    \underline{Induction step:} Suppose the result is true for any simple orthogeodesic which intersects $\tau_1$ at most $n$ times before intersecting $\tau_2$. Let $\tau$ be a simple orthogeodesic which intersects $\tau_1$ $n+1$ times before intersecting $\tau_2$.\\
    
    We denote by $p_1, p_2, ..., p_n, p_{n+1}$ the first $n+1$ intersection points of $\tau$ and $\tau_1$. These points cut $\tau_1$ into $n+2$ segments.
    
    \begin{figure}[H]
        \centering
        \includegraphics[height=12.5cm]{figures/OctogoneTauCas2Etape1.png}
        \caption{Several example of $\tau$.}
        \label{fig:octoCas2Step1}
    \end{figure}
    
    Let us construct two arcs $\sigma$ and $\Tilde{\sigma}$ as follows. For~$\sigma$, we start from the edge~$b_2$ and we follow~$\tau$ until it intersects~$\tau_1$, then we follow~$\tau_1$ until~$p_n$, then we follow~$\tau$ between~$p_n$ and~$p_{n-1}$, then~$\tau_1$ between~$p_{n-1}$ and~$p_{n-2}$ and so on until we reach~$p_1$. If~$n+1$ is even, we close up~$\sigma$ by following~$\tau$ until the edge~$b_1$; if~$n+1$ is odd, we follow~$\tau_1$ until the edge~$b_1$.
    We construct~$\Tilde{\sigma}$ in a similar way, using segments of~$\tau$ and~$\tau_1$ between the intersection points~$p_i$ that we did not already use. We start from the side~$b_2$ and follow~$\tau_1$ until~$p_{n+1}$, then we follow~$\tau$ until~$p_n$, then~$\tau_1$ until~$p_{n-1}$ and we repeat the process until we reach~$p_1$. Then, if~$n+1$ is even, we follow~$\tau_1$ until~$b_1$; if~$n+1$ is odd, we follow~$\tau$ until~$b_1$.

    \begin{figure}[H]
        \centering
        \includegraphics[height=12.5cm]{figures/OctogoneTauCas2Etape2.png}
        \caption{Several examples of the construction of $\sigma$ and $\Tilde{\sigma}$ in Case 2.}
        \label{fig:octoCas2Step2}
    \end{figure}
    Since $\sigma$ and $\Tilde{\sigma}$ use different segments of $\tau$ and $\tau_1$, we have $\ell(\sigma) + \ell(\Tilde{\sigma}) \leqslant \ell(\tau_1) + \ell(\tau) $. By induction, we know that $\ell(\tau_3) < \ell(\sigma)$ and $\ell(\tau_3) < \ell(\Tilde{\sigma})$. Indeed, when $n+1$ is even, the arc $\sigma$ is homotopic to $\tau$ in the induction step $\frac{n+1}{2}$ for $i(\tau,\tau_2)=0$ and the arc $\Tilde{\sigma}$ is homotopic to $\tau$ in the induction step $(\frac{n+1}{2}-1)$ for $i(\tau,\tau_2)=0$. When $n+1$ is odd, the arcs $\sigma$ and $\Tilde{\sigma}$ are homotopic to $\tau$ in the induction step $\frac{n}{2}$ for $i(\tau,\tau_2)=0$. To conclude, we have $2\ell(\tau_3) < \ell(\sigma) + \ell(\Tilde{\sigma}) \leqslant \ell(\tau) + \ell(\tau_1)$. Thus, $\ell(\tau_3) < \ell(\tau)$.
    \\
    
    So the first three lengths of $\mathcal{O}_S(T)$ and $\mathcal{O}_S(T')$ are realized by disjoint orthogeodesics~$\tau_1, \tau_2$ and $\tau_3$ on $T$ and $\tau_1', \tau_2'$ and $\tau_3'$ on $T'$. The set $\{ \tau_1$, $\tau_2, \tau_3 \}$ is a hexagon decomposition of $T$ and the set $\{\tau'_1$, $\tau'_2, \tau'_3\}$ is also a hexagon decomposition of~$T'$. Cutting $T$ and $T'$ along their respective hexagon decomposition, we obtain two isometric sets of two hexagons. We have only two hexagons per set and they are isometric, so there is no ambiguity as to how to glue them back into $T$ and $T'$, which are then isometric.  
\end{proof}



Finally, in~\cite{orthoSys}, Masai and McShane also gave an example of two non-isometric hyperbolic surfaces with the same orthospectrum. Their proof does not provide such an example in the case of the simple orthospectrum. Indeed, they used the fact that if we have a regular~$d$-cover $\pi : \Tilde{X} \to X$ of a hyperbolic surface $X$ with boundary, then any orthogeodesic on $X$ is covered by exactly $d$ orthogeodesics on $\Tilde{X}$ \cite[Lemma~6.1]{orthoSys}. It is then possible to compute the orthospectrum of $\Tilde{X}$ from $\mathcal{O}(X)$ and the degree of the cover. They construct two non-isometric regular degree $d$ cover of the same hyperbolic surface, which then have the same orthospectrum. To use the same argument for the simple orthospectrum, we would need to control which non-simple orthogeodesics on $X$ have simple lift to $\Tilde{X}$. So, similarly to the simple spectrum case, the question of whether the simple orthospectrum determine the surface is still open.

\section{Conflict of interest statement.}
There is no conflict of interest. Data sharing not applicable to this
article as no datasets were generated or analyzed during the current study.


\bibliographystyle{alpha}
\bibliography{bibliography}

\begin{flushright}
    \address \\
    Email address: \email
\end{flushright}
    
%%% Merci pour votre lecture !

\end{document}
